% Luc Destombe --- licence CC BY-NC-SA 4.0
% https://creativecommons.org/licenses/by-nc-sa/4.0/deed.fr
% Fichier autonome : compiler avec pdflatex, deux fois (signets, nombre de pages).
\def\lyceeversion{1.1}
\documentclass[11pt,a4paper]{article}
\makeatletter
\newif\ifeleve
\elevefalse

\newif\iflycee@palatino
\lycee@palatinotrue

\RequirePackage[utf8]{inputenc}
\RequirePackage[T1]{fontenc}
\RequirePackage[french]{babel}
\RequirePackage{etoolbox}

\newcommand{\ShorthandsOff}{\@ifundefined{shorthandoff}{}{\shorthandoff{;:!?}}}
\newcommand{\ShorthandsOn}{\@ifundefined{shorthandon}{}{\shorthandon{;:!?}}}
\ShorthandsOff
\AtBeginDocument{\ShorthandsOn}
\AtBeginEnvironment{tikzpicture}{\ShorthandsOff}

\RequirePackage{geometry}
\RequirePackage{fancyhdr}
\RequirePackage{lastpage}
\RequirePackage{multicol}
\RequirePackage{array}
\RequirePackage{enumitem}
\RequirePackage{xcolor}
\RequirePackage{needspace}

\RequirePackage{amsmath,amssymb}
\RequirePackage{mathpazo}
\renewcommand{\ttdefault}{lmtt}
\DeclareMathSizes{10.95}{10.95}{8}{5.5}
\begingroup\catcode`\_=\active \catcode`\^=\active
\gdef_#1{\sb{\scriptscriptstyle #1}}
\gdef^#1{\sp{\scriptscriptstyle\lycee@moinsepais #1}}
\endgroup
\newcommand\lycee@moins{\mathbin{%
  \setbox\z@\hbox{$\m@th\scriptscriptstyle\lycee@moinsorig$}%
  \dimen@=.55\wd\z@ \advance\dimen@-.5pt
  \hbox to.75\wd\z@{\hss$\m@th\scriptscriptstyle\vcenter{\hbox{%
    \tikz\draw[line width=.5pt,line cap=round](0,0)--(\the\dimen@,0);}}$\hss}}}
\begingroup\lccode`\~=`\- \lowercase{\endgroup\let~}\lycee@moins
\newcommand\lycee@moinsepais{\mathcode`\-="8000 }
\AtBeginDocument{\mathchardef\lycee@moinsorig=\mathcode`\-\relax}
\AtBeginDocument{\catcode`\_=12 \mathcode`\_="8000
  \catcode`\^=12 \mathcode`\^="8000 }
\newcommand\lycee@exposants{%
  \fontdimen13\textfont2=.48\fontdimen6\textfont2
  \fontdimen14\textfont2=.48\fontdimen6\textfont2
  \fontdimen15\textfont2=.48\fontdimen6\textfont2
  \fontdimen13\scriptfont2=.48\fontdimen6\scriptfont2
  \fontdimen14\scriptfont2=.48\fontdimen6\scriptfont2
  \fontdimen15\scriptfont2=.48\fontdimen6\scriptfont2}
\every@math@size\expandafter{\the\every@math@size\lycee@exposants}
\newlength{\retraitcalculs}
\setlength{\retraitcalculs}{2em}

\RequirePackage{tikz}
\usetikzlibrary{arrows.meta,calc,babel,shapes.misc}
\RequirePackage[most]{tcolorbox}
\tcbuselibrary{skins,breakable}

\definecolor{coulDef}{HTML}{1F4E79}
\definecolor{coulProp}{HTML}{7030A0}
\definecolor{coulRem}{HTML}{C55A11}
\definecolor{coulEx}{HTML}{2E7D32}
\definecolor{coulExo}{HTML}{B03A2E}
\definecolor{coulPart}{HTML}{1F4E79}
\definecolor{coulMeth}{HTML}{1B6E4F}
\definecolor{coulCourbe}{HTML}{0B5ED7}
\definecolor{coulCourbeB}{HTML}{C00000}
\definecolor{coulCourbeC}{HTML}{2E7D32}
\definecolor{coulGrille}{HTML}{8C8C8C}
\definecolor{coulRep}{HTML}{1B6E4F}
\definecolor{coulBar}{HTML}{2E7D32}

\newcommand{\R}{\mathbb{R}}
\newcommand{\Rs}{\mathbb{R}^{*}}
\renewcommand{\le}{\leqslant}
\renewcommand{\ge}{\geqslant}
\renewcommand{\approx}{\simeq}
\newcommand{\vect}[1]{\overrightarrow{#1}}
\newcommand{\norme}[1]{\left\lVert #1 \right\rVert}
\newcommand{\intervalle}[2]{\left[#1\,;#2\right]}
\RequirePackage{cancel}
\renewcommand{\CancelColor}{\color{coulExo}}
\newcommand{\fois}[1]{\times{\color{coulExo}#1}}
\newcommand{\barre}[1]{\cancel{#1}}

\tikzset{grille/.style={coulGrille,line width=0.4pt}}
\newcommand{\quadrillage}[4]{\draw[grille,step=1] (#1,#2) grid (#3,#4);}
\tikzset{
  croix/.style={cross out,draw,line width=0.6pt,inner sep=0pt,minimum size=4pt},
  croixdroite/.style={croix,rotate=45,line width=0.8pt},
}
\newcommand{\pt}[3]{\path (#1) node[croix]{} node[#2,font=\small,inner sep=2.5pt]{$#3$};}

\newcommand{\lycee@repere}[4]{%
  \draw[grille,step=1] (#1,#2) grid (#3,#4);
  \draw[-{Stealth[length=2mm]},thick] (#1,0)--({#3+0.4},0) node[below left=-1pt]{$x$};
  \draw[-{Stealth[length=2mm]},thick] (0,#2)--(0,{#4+0.4}) node[below left=-1pt]{$y$};
  \node[below left,font=\scriptsize,inner sep=1.5pt] at (0,0) {$0$};}
\newcommand{\lycee@gradx}[1]{\foreach \k/\l in {#1}{%
  \draw (\k,0.07)--(\k,-0.07) node[below,font=\scriptsize,inner sep=1.5pt]{$\l$};}}
\newcommand{\lycee@grady}[1]{\foreach \k/\l in {#1}{%
  \draw (0.07,\k)--(-0.07,\k) node[left,font=\scriptsize,inner sep=1.5pt]{$\l$};}}
\newcommand{\lycee@intff}[2]{\left[#1\,;#2\right]}
\newcommand{\lycee@intfo}[2]{\left[#1\,;#2\right[}
\newcommand{\lycee@intof}[2]{\left]#1\,;#2\right]}
\newcommand{\lycee@intoo}[2]{\left]#1\,;#2\right[}
\newcommand{\lycee@ens}[1]{\left\{#1\right\}}
\AtBeginDocument{%
  \providecommand{\repere}{\lycee@repere}%
  \providecommand{\gradx}{\lycee@gradx}%
  \providecommand{\grady}{\lycee@grady}%
  \providecommand{\Cf}{\mathcal{C}\sb{\scriptscriptstyle f}}%
  \providecommand{\Cg}{\mathcal{C}\sb{\scriptscriptstyle g}}%
  \providecommand{\intff}{\lycee@intff}%
  \providecommand{\intfo}{\lycee@intfo}%
  \providecommand{\intof}{\lycee@intof}%
  \providecommand{\intoo}{\lycee@intoo}%
  \providecommand{\ens}{\lycee@ens}}

\newlist{questions}{enumerate}{3}
\setlist[questions,1]{label=\arabic*\textdegree),leftmargin=*,itemsep=2pt,topsep=3pt}
\setlist[questions,2]{label=\alph*),leftmargin=*,itemsep=1pt,topsep=2pt}
\setlist[questions,3]{label=\roman*),leftmargin=*,itemsep=1pt,topsep=1pt}
\setlist[itemize]{label=\textbullet,leftmargin=*,itemsep=2pt,topsep=3pt}

\newcommand{\besoinplace}[1]{\par\penalty\@M\begingroup
  \dimen@=#1\relax\dimen@ii=\pagegoal\advance\dimen@ii-\pagetotal
  \ifdim\dimen@>\dimen@ii\ifdim\dimen@ii>\z@\vfil\fi\break\fi\endgroup}

\newcommand{\lycee@pied}{\thepage/\pageref{LastPage}}
\newcommand{\entete}[2]{%
  \pagestyle{fancy}\fancyhf{}%
  \fancyhead[L]{\footnotesize\color{coulPart}\textbf{#1}}%
  \fancyhead[R]{\footnotesize\color{coulPart}\textbf{#2}}%
  \fancyfoot[C]{\footnotesize\lycee@pied}%
  \renewcommand{\headrulewidth}{0.4pt}%
  \renewcommand{\headrule}{\hbox to\headwidth{%
    \color{coulPart!40}\leaders\hrule height \headrulewidth\hfill}}}

\RequirePackage{ccicons}
\newcommand{\lycee@licence}{{\scriptsize\color{gray}%
  \copyright\ Luc Destombe --- \ccbyncsa\ CC BY-NC-SA 4.0}}
\newif\iflycee@licence \lycee@licencetrue
\newcommand{\sanslicence}{\lycee@licencefalse}
\AtBeginDocument{\iflycee@licence\AddToHookNext{shipout/foreground}{%
  \put(0,0){\rlap{\kern\dimexpr1in+\hoffset+\oddsidemargin+\textwidth\relax
    \llap{\raisebox{-\dimexpr1in+\voffset+\topmargin+\headheight+\headsep
      +\textheight+\footskip\relax}{\lycee@licence}}}}}\fi}

\newcommand{\lycee@site}{faire-des-maths.fr}
\newcommand{\lycee@version}{{\scriptsize\color{gray}\lycee@site
  \ifdefined\lyceeversion\ --- Version~\lyceeversion\fi}}
\AtBeginDocument{\iflycee@licence\AddToHookNext{shipout/foreground}{%
  \put(0,0){\rlap{\kern\dimexpr1in+\hoffset+\oddsidemargin\relax
    \raisebox{-\dimexpr1in+\voffset+\topmargin+\headheight+\headsep
      +\textheight+\footskip\relax}{\lycee@version}}}}\fi}

\newcommand{\formulesserrees}{%
  \setlength{\abovedisplayskip}{3pt}\setlength{\belowdisplayskip}{3pt}%
  \setlength{\abovedisplayshortskip}{2pt}\setlength{\belowdisplayshortskip}{3pt}}

\setlength{\parindent}{0pt}

\RequirePackage[implicit=false,hidelinks,pdfencoding=auto,psdextra,bookmarksopen,
  bookmarksopenlevel=1,pdfcreator={},pdfproducer={}]{hyperref}
\RequirePackage{bookmark}
\hypersetup{pdfauthor={Luc Destombe},pdfkeywords={CC BY-NC-SA 4.0}}
\newcounter{lycee@signet}
\newcommand{\signet}[3][]{\stepcounter{lycee@signet}%
  \edef\lycee@dest{\if\relax\detokenize{#1}\relax
    signet.\arabic{lycee@signet}\else#1\fi}%
  \hypertarget{\lycee@dest}{}\bookmark[level=#2,dest=\lycee@dest]{#3}}
\pdfstringdefDisableCommands{%
  \def\R{R}\def\vect#1{#1}\def\Cf{Cf}\def\Cg{Cg}\def\fois#1{×#1}%
  \def\barre#1{#1}\def\intervalle#1#2{[#1;#2]}\def\intff#1#2{[#1;#2]}%
  \def\textsuperscript#1{#1}\def\dfrac#1#2{#1/#2}\def\frac#1#2{#1/#2}%
  \def\vec#1{#1⃗}}%
\newcommand{\res}[1]{#1}
\geometry{top=2cm,bottom=1cm,left=1cm,right=1cm,headheight=14pt,footskip=0.6cm}

\RequirePackage{pgfplots}
\pgfplotsset{compat=1.18}
\RequirePackage{tkz-tab}

\newcounter{partiecours}
\renewcommand{\thepartiecours}{\Roman{partiecours}}
\newcommand{\partiecours}[1]{%
  \refstepcounter{partiecours}%
  \Needspace*{8\baselineskip}%
  \bigskip
  \noindent\begin{tcolorbox}[enhanced,colback=coulPart,colframe=coulPart,
      boxrule=0pt,arc=2pt,left=8pt,right=8pt,top=4pt,bottom=4pt,
      before skip=6pt,after skip=10pt]
    \color{white}\bfseries\large \thepartiecours{} --- #1%
    \signet[partie-\arabic{partiecours}]{1}{\thepartiecours{} --- #1}
  \end{tcolorbox}%
  \addcontentsline{toc}{section}{\thepartiecours{} --- #1}%
}

\newtcolorbox{boitecours}[3][]{%
  enhanced,breakable,
  colback=white,colframe=#2,coltitle=white,
  fonttitle=\bfseries,
  attach boxed title to top left={xshift=8pt,yshift=-\tcboxedtitleheight/2},
  boxed title style={colback=#2,arc=2pt,boxrule=0pt},
  boxrule=0.9pt,arc=2pt,
  left=8pt,right=8pt,top=8pt,bottom=6pt,
  title={#3},#1}

\newcounter{methode}
\newenvironment{definitionbox}[1][Définition]%
  {\begin{boitecours}[phantom={\signet{2}{#1}}]{coulDef}{#1}}{\end{boitecours}}
\newenvironment{proprietebox}[1][Propriété]%
  {\begin{boitecours}[phantom={\signet{2}{#1}}]{coulProp}{#1}}{\end{boitecours}}
\newenvironment{methodebox}[1][Méthode]{\stepcounter{methode}%
  \begin{boitecours}[phantom={\signet[methode-\arabic{methode}]{2}{#1}}]{coulMeth}{#1}}%
  {\end{boitecours}}

\newcommand{\etape}[1]{\par\smallskip\textbf{\color{coulMeth}#1}\par\nobreak\smallskip}
\newcommand{\enonce}[1]{\emph{#1}\par\smallskip}
\newcommand{\reponse}[1]{\par\smallskip
  {\footnotesize\itshape\color{black!55}Réponses : #1}}
\newcommand{\demo}[1]{\textbf{\color{coulMeth}Démonstration.} #1}
\newcommand{\afaire}[1]{%
  \par\smallskip
  \noindent{\small\color{coulExo}$\blacktriangleright$\ \textbf{À faire :}
    fiche d'exercices du chapitre, #1.}\par\smallskip}

\newenvironment{calculs}{\setlength{\jot}{5pt}\@fleqntrue\@mathmargin\retraitcalculs
  \setlength{\abovedisplayskip}{4pt}\setlength{\belowdisplayskip}{4pt}%
  \setlength{\abovedisplayshortskip}{4pt}\setlength{\belowdisplayshortskip}{4pt}%
  \csname alignat*\endcsname{2}}%
  {\csname endalignat*\endcsname}
\newcommand{\rem}[1]{\text{\small\itshape\color{black!60}#1}}
\newcommand{\explique}[2]{\par\smallskip\noindent
  \begin{minipage}[t]{0.57\linewidth}\raggedright #1\end{minipage}\hfill
  \begin{minipage}[t]{0.40\linewidth}\raggedright\leavevmode
    {\small\itshape\color{black!60}#2}\end{minipage}\par}
\newcommand{\cas}[1]{\par\smallskip\noindent\textbf{\color{coulExo}#1}\vadjust{\nobreak}}
\newcommand{\calc}[1]{\par\smallskip\textbf{\color{coulExo}#1}\par\nobreak}

\newtcolorbox{remarquebox}[1][Remarque]{%
  enhanced,breakable,colback=white,colframe=coulRem,
  boxrule=0pt,leftrule=2.5pt,arc=0pt,outer arc=0pt,
  left=8pt,right=6pt,top=5pt,bottom=5pt,
  before upper={\textbf{\color{coulRem}#1}\par\smallskip}}

\newtcolorbox{exemplebox}[1][Exemple]{%
  enhanced,breakable,colback=white,colframe=coulEx,
  boxrule=0pt,leftrule=2.5pt,arc=0pt,outer arc=0pt,
  left=8pt,right=6pt,top=5pt,bottom=5pt,
  before upper={\textbf{\color{coulEx}#1}\par\smallskip}}

\pgfplotsset{
  repere/.style={
    axis lines=middle,
    axis line style={-{Stealth[length=2.4mm]},thick},
    every axis x label/.style={at={(ticklabel* cs:1.0)},anchor=west},
    every axis y label/.style={at={(ticklabel* cs:1.0)},anchor=south},
    label style={font=\footnotesize},
    tick label style={font=\scriptsize},
    grid=both,
    major grid style={grille},
    minor tick num=0,
    tick align=inside,
    clip mode=individual,
  },
}
\tikzset{
  courbe/.style={coulCourbe,line width=1.1pt,smooth},
  courbeB/.style={coulCourbeB,line width=1.1pt,smooth},
  courbeC/.style={coulCourbeC,line width=1.1pt,smooth},
}

\newlist{etapes}{enumerate}{1}
\setlist[etapes]{label=\textbf{\arabic*.},leftmargin=*,itemsep=1pt,topsep=2pt}

\newlist{expressions}{enumerate}{1}
\setlist[expressions]{label={\Alph*\ $=$},leftmargin=*,itemsep=3pt,topsep=3pt}

\newcounter{exo}
\renewcommand{\theexo}{\ifnum\value{exo}<10 0\fi\arabic{exo}}
\newtcolorbox{exercice}[1][]{%
  enhanced,breakable,
  colback=white,colframe=coulExo,
  boxrule=0.9pt,arc=2pt,
  left=8pt,right=8pt,top=8pt,bottom=6pt,
  coltitle=white,fonttitle=\bfseries,
  attach boxed title to top left={xshift=8pt,yshift=-\tcboxedtitleheight/2},
  boxed title style={colback=coulExo,arc=2pt,boxrule=0pt},
  phantom={\signet{2}{Exercice \arabic{exo}}},
  title={Exercice \theexo\ifx\relax#1\relax\else{} --- #1\fi}}
\newenvironment{exo}[1][\relax]{\refstepcounter{exo}\begin{exercice}[#1]}{\end{exercice}}

  \newenvironment{correction}{\par}{\par}
  \newcommand{\autableau}{}
  \newcommand{\etapeexemples}{\etape{Exemples corrigés}}
  \newcommand{\etapetableau}[1]{\etape{#1}}
  \newenvironment{exempletableau}[1][Exemple]%
    {\begin{exemplebox}[#1]}{\end{exemplebox}}

\RequirePackage{comment}
\excludecomment{correctionavj}

\newcommand{\coursEntete}{}
\newcommand{\coursNiveau}{}
\newcommand{\coursTitre}{}
\newcommand{\coursSurTitre}{}
\newcommand{\coursSousTitre}{}

\newcommand{\enteteCours}[1]{\renewcommand{\coursEntete}{#1}}
\newcommand{\chapitrecours}[2]{\enteteCours{Chapitre #1 --- #2}}
\newcommand{\niveaucours}[1]{\renewcommand{\coursNiveau}{#1}}
\newcommand{\titrecours}[3][]{%
  \renewcommand{\coursTitre}{#2}%
  \renewcommand{\coursSurTitre}{#1}%
  \renewcommand{\coursSousTitre}{#3}}

\pagestyle{fancy}
\fancyhf{}
\fancyhead[L]{\footnotesize\color{coulPart}\textbf{\coursEntete}}
\fancyhead[R]{\footnotesize\color{coulPart}\textbf{\coursNiveau}}
\fancyfoot[C]{\footnotesize\thepage}
\renewcommand{\headrulewidth}{0.4pt}
\renewcommand{\headrule}{\hbox to\headwidth{\color{coulPart!40}\leaders\hrule height \headrulewidth\hfill}}

\setlength{\parskip}{2pt}

\newcommand{\lycee@mentioneleve}{\\[2pt]{\small\itshape Version élève}}
\AtBeginDocument{%
  \begin{center}
    \begin{tcolorbox}[enhanced,width=0.72\linewidth,colback=white,colframe=coulPart,
        boxrule=1.6pt,arc=3pt,halign=center,top=10pt,bottom=10pt]
      \ifx\coursSurTitre\empty
        {\Huge\bfseries\color{coulPart} \coursTitre}\\[4pt]
      \else
        {\Huge\bfseries\color{coulPart} \coursTitre}\\[3pt]
        {\Large\bfseries\color{coulPart} \coursSurTitre}\\[5pt]
      \fi
      {\normalsize \coursSousTitre}%
      
    \end{tcolorbox}
  \end{center}%
}

\RequirePackage{tikz}
\RequirePackage{xcolor}

\usetikzlibrary{decorations.text}
\definecolor{lycee@metal}{HTML}{1F4E79}
\definecolor{lycee@manche}{HTML}{8A5A2B}
\definecolor{lycee@etincelle}{HTML}{F28C28}
\newcommand{\lycee@pluriel}[2]{#1 #2\ifnum#1>1 s\fi}
\newcommand{\lycee@sceau}{%
  \edef\lycee@bas{\lycee@pluriel{\lyceerelectures}{relecture} et
    \lycee@pluriel{\lyceerodages}{rodage}}%
  \begin{tikzpicture}[line join=round]
    \fill[coulPart!10,draw=coulPart,line width=.8pt] (0:2.2)
      \foreach \i in {1,...,71}{ -- (\i*5:{isodd(\i)?2.1:2.2}) } -- cycle;
    \draw[coulPart,line width=1.2pt] (0,0) circle (1.93);
    \draw[coulPart,line width=.4pt] (0,0) circle (1.85);
    \draw[coulPart,line width=.4pt] (0,0) circle (1.2);
    \path[decorate,decoration={text along path,text align=center,
        text={|\fontsize{9.5}{11}\selectfont\bfseries\color{coulPart}|TEST{\'E} AVEC MES {\'E}L{\`E}VES}}]
      (180:1.45) arc (180:0:1.45);
    \path[decorate,decoration={text along path,text align=center,
        text/.expanded={|\noexpand\fontsize{8.5}{10}\noexpand\selectfont\noexpand\color{coulPart}|\lycee@bas}}]
      (195:1.68) arc (195:345:1.68);
    \begin{scope}[x=.25mm,y=-.25mm]
      \begin{scope}[shift={(-32,-40)}]%
        \fill[lycee@metal] (23,31) rectangle (41,49);
        \draw[lycee@etincelle,line width=.5pt,line cap=round]
          (20,30)--(15,27) (19,34)--(14,35) (44,30)--(49,27) (45,34)--(50,35);
        \fill[lycee@manche,rounded corners=.3mm] (38,24.6) rectangle (60,27.6);
        \fill[lycee@metal,rounded corners=.4mm] (25,20) rectangle (39,31);
        \fill[lycee@etincelle,rounded corners=.2mm] (25,20) rectangle (28,31);
      \end{scope}
    \end{scope}
    \ifdefined\lyceesceauplusieurs\def\lycee@annee{Ann{\'e}es~:}\else
      \def\lycee@annee{Ann{\'e}e}\fi
    \path[decorate,decoration={text along path,text align=center,
        text/.expanded={|\noexpand\fontsize{8}{10}\noexpand\selectfont\noexpand\color{coulPart}|%
          \unexpanded\expandafter{\lycee@annee} \unexpanded\expandafter{\lyceesceauannees}}}]
      (95:1.08) arc (95:445:1.08);
  \end{tikzpicture}}
\AtEndDocument{\ifdefined\lyceesceauannees
  \par\besoinplace{5cm}\vspace{1.5em}%
  \begin{center}\lycee@sceau\end{center}\fi}
\makeatother

\chapitrecours{9}{Géométrie repérée et produit scalaire}
\niveaucours{1\textsuperscript{re} spécialité mathématiques}
\titrecours{GÉOMÉTRIE REPÉRÉE ET PRODUIT SCALAIRE}{Cours --- Première spécialité}

\begin{document}

\newcommand{\coord}[2]{\begin{pmatrix}#1\\#2\end{pmatrix}}

\ShorthandsOff
\tikzset{
  vecteur/.style={-{Stealth[length=2.6mm]},line width=1.1pt,coulCourbe},
  vecteurB/.style={-{Stealth[length=2.6mm]},line width=1.1pt,coulCourbeB},
  vecteurC/.style={-{Stealth[length=2.6mm]},line width=1.1pt,coulCourbeC},
  base/.style={-{Stealth[length=2mm]},line width=0.9pt,black},
  aide/.style={dashed,black!45,line width=0.6pt},
  fig/.style={line width=0.8pt,black!75},
  axe/.style={aide,-{Stealth[length=2.2mm]}},
}
\ShorthandsOn

\newcommand{\carremilieu}[5]{%
\begin{tikzpicture}[scale=0.45]
  \draw[fig] (0,0) -- (4,0) -- (4,4) -- (0,4) -- cycle;
  \draw[aide] (0,0) -- (4,4);
  \pt{0,0}{below left}{#1} \pt{4,0}{below right}{#2} \pt{4,4}{above right}{#3}
  \pt{0,4}{above left}{#4} \pt{2,4}{above}{#5}
\end{tikzpicture}}

\begin{remarquebox}[Comment utiliser ce cours]
Chaque partie commence par ce qu'il faut \textbf{savoir} (définitions, théorèmes), puis
vient ce qu'il faut \textbf{savoir faire} : une méthode, avec des
\textbf{exemples corrigés} faits en classe, puis des exercices
\og\textbf{À vous de jouer}\fg{} à chercher seul.

\end{remarquebox}

\partiecours{Coordonnées d'un vecteur, somme de deux vecteurs}

Les parties I à IV reprennent la seconde (chapitres « Vecteurs » et « Vecteurs et
trigonométrie »). Le plan est muni d'un repère $(O\,;\vec{\imath},\vec{\jmath})$ ; à partir
de la partie IV, il est \textbf{orthonormé}.

\begin{definitionbox}[Définition 1 --- Coordonnées d'un vecteur]
\begin{minipage}[t]{0.58\linewidth}\raggedright
Tout vecteur $\vec{u}$ s'écrit d'une seule façon $\vec{u}=x\,\vec{\imath}+y\,\vec{\jmath}$,
avec deux nombres réels $x$ et $y$ : ce sont les \textbf{coordonnées} de $\vec{u}$. On
note $\vec{u}\coord{x}{y}$.

\smallskip
Sur la figure, $\vec{u}=3\vec{\imath}+2\vec{\jmath}$ : on avance de $3$ vers la droite et
de $2$ vers le haut, donc $\vec{u}\coord{3}{2}$.
\end{minipage}\hfill
\begin{minipage}[t]{0.38\linewidth}
\vspace{0pt}\centering
\begin{tikzpicture}[x=0.6cm,y=0.6cm]
  \quadrillage{-0.4}{-0.4}{4.4}{3.4}
  \draw[base] (0,0) -- node[below,font=\footnotesize]{$\vec{\imath}$} (1,0);
  \draw[base] (0,0) -- node[left,font=\footnotesize]{$\vec{\jmath}$} (0,1);
  \draw[aide] (1,1) -- (4,1) -- (4,3);
  \draw[vecteur] (1,1) -- node[midway,above left,font=\footnotesize]{$\vec{u}$} (4,3);
  \node[below,font=\footnotesize] at (2.5,1) {$3$};
  \node[right,font=\footnotesize] at (4,2) {$2$};
  \pt{0,0}{below left}{O}
\end{tikzpicture}
\end{minipage}
\end{definitionbox}

\begin{proprietebox}[Propriété 1 --- Somme et différence]
Pour deux vecteurs $\vec{u}\coord{x}{y}$ et $\vec{v}\coord{x'}{y'}$ :
\begin{center}
  $\vec{u}+\vec{v}\coord{x+x'}{y+y'}$ \qquad $\vec{u}-\vec{v}\coord{x-x'}{y-y'}$
\end{center}
On ajoute (ou on soustrait) les abscisses entre elles et les ordonnées entre elles.
\end{proprietebox}

\begin{proprietebox}[Propriété 2 --- Coordonnées de $\vect{AB}$]
Pour deux points $A(x_{A}\,;y_{A})$ et $B(x_{B}\,;y_{B})$ :
\[
  \vect{AB}\coord{x_{B}-x_{A}}{y_{B}-y_{A}}
  \qquad\text{(arrivée moins départ).}
\]
\demo{par la relation de Chasles :}
\begin{calculs}
\vect{AB} &= \vect{AO}+\vect{OB} &\qquad& \rem{relation de Chasles}\\
\vect{AB} &= \vect{OB}-\vect{OA} && \rem{$\vect{AO}=-\vect{OA}$}
\end{calculs}
Or $\vect{OB}$ a pour coordonnées celles du point $B$, et $\vect{OA}$ celles de $A$ : la
propriété 1 donne le résultat.
\end{proprietebox}

\begin{methodebox}[Méthode 1 --- Calculer les coordonnées d'un vecteur]
\etapeexemples
\begin{questions}[label=\alph*)]
  \item Donner les coordonnées des vecteurs $\vec{w}=4\vec{\imath}-\vec{\jmath}$ et
        $\vec{t}=-2\vec{\jmath}$.
  \item Pour les points $A(3\,;-1)$ et $B(-2\,;4)$, calculer les coordonnées de
        $\vect{AB}$, puis celles de $\vect{BA}$.
  \item Pour les points $E\left(\frac{2}{3}\,;1\right)$ et $F\left(-1\,;\frac{5}{3}\right)$,
        calculer les coordonnées de $\vect{EF}$.
\end{questions}
\begin{correction}
\cas{a)}
\explique{$\vec{w}=4\vec{\imath}+(-1)\vec{\jmath}$, donc $\res{\vec{w}\coord{4}{-1}}$.}{le
nombre devant $\vec{\imath}$, puis celui devant $\vec{\jmath}$}
\explique{$\vec{t}=0\,\vec{\imath}+(-2)\vec{\jmath}$, donc $\res{\vec{t}\coord{0}{-2}}$.}{pas
de $\vec{\imath}$ : l'abscisse est nulle}
\cas{b)}
\begin{calculs}
\vect{AB} & \coord{x_{B}-x_{A}}{y_{B}-y_{A}} &\qquad& \rem{arrivée moins départ}\\
\vect{AB} & \coord{-2-3}{4-(-1)} && \rem{on remplace par les coordonnées}\\
\vect{AB} & \res{\coord{-5}{5}} &&\\
\vect{BA} & \coord{3-(-2)}{-1-4} && \rem{arrivée $A$, départ $B$}\\
\vect{BA} & \res{\coord{5}{-5}} && \rem{l'opposé de $\vect{AB}$ : le trajet inverse}
\end{calculs}
\cas{c)}
\begin{calculs}
x_{F}-x_{E} &= -1-\dfrac{2}{3} &\qquad& \rem{l'abscisse de $\vect{EF}$}\\
x_{F}-x_{E} &= -\dfrac{1\fois{3}}{1\fois{3}}-\dfrac{2}{3} && \rem{même dénominateur}\\
x_{F}-x_{E} &= -\dfrac{5}{3} &&\\
y_{F}-y_{E} &= \dfrac{5}{3}-1 && \rem{l'ordonnée de $\vect{EF}$}\\
y_{F}-y_{E} &= \dfrac{5}{3}-\dfrac{1\fois{3}}{1\fois{3}} &&\\
y_{F}-y_{E} &= \dfrac{2}{3} &&
\end{calculs}
\explique{Le vecteur $\vect{EF}$ a pour coordonnées
$\res{\vect{EF}\coord{-\frac{5}{3}}{\frac{2}{3}}}$.}{on conclut}
\end{correction}

\etape{À vous de jouer}
\begin{questions}[label=\alph*)]
  \item Donner les coordonnées des vecteurs $\vec{w}=-3\vec{\imath}+5\vec{\jmath}$ et
        $\vec{s}=6\vec{\imath}$.
  \item Pour les points $C(-4\,;2)$ et $D(1\,;-1)$, calculer les coordonnées de
        $\vect{CD}$, puis celles de $\vect{DC}$.
  \item Pour les points $G\left(\frac{1}{3}\,;-2\right)$ et $H\left(2\,;-\frac{4}{3}\right)$,
        calculer les coordonnées de $\vect{GH}$.
\end{questions}
\reponse{a)~$\vec{w}\coord{-3}{5}$ et $\vec{s}\coord{6}{0}$ ;\; b)~$\vect{CD}\coord{5}{-3}$ et
$\vect{DC}\coord{-5}{3}$ ;\; c)~$\vect{GH}\coord{\frac{5}{3}}{\frac{2}{3}}$.}
\end{methodebox}

\begin{methodebox}[Méthode 2 --- Calculer avec les coordonnées]
\etapeexemples
On considère les vecteurs $\vec{u}\coord{3}{-2}$ et $\vec{v}\coord{-1}{5}$, et les points
$A(2\,;1)$, $B(-1\,;3)$ et $C(4\,;-2)$.
\begin{questions}[label=\alph*)]
  \item Calculer les coordonnées des vecteurs $\vec{u}+\vec{v}$ et $\vec{u}-\vec{v}$.
  \item Calculer les coordonnées du vecteur $\vect{AB}+\vect{BC}$. Quel vecteur
        retrouve-t-on ?
  \item Calculer les coordonnées du vecteur $\vec{w}$ tel que $\vec{u}+\vec{w}=\vec{v}$.
\end{questions}
\begin{correction}
\cas{a)}
\begin{calculs}
\vec{u}+\vec{v} & \coord{3+(-1)}{-2+5} &\qquad& \rem{abscisses ensemble, ordonnées ensemble}\\
\vec{u}+\vec{v} & \res{\coord{2}{3}} &&\\
\vec{u}-\vec{v} & \coord{3-(-1)}{-2-5} && \rem{attention aux signes}\\
\vec{u}-\vec{v} & \res{\coord{4}{-7}} &&
\end{calculs}
\cas{b)}
\begin{calculs}
\vect{AB} & \coord{-1-2}{3-1} &\qquad& \rem{d'abord chaque vecteur : arrivée moins départ}\\
\vect{AB} & \coord{-3}{2} &&\\
\vect{BC} & \coord{4-(-1)}{-2-3} &&\\
\vect{BC} & \coord{5}{-5} &&\\
\vect{AB}+\vect{BC} & \coord{-3+5}{2+(-5)} && \rem{puis la somme}\\
\vect{AB}+\vect{BC} & \res{\coord{2}{-3}} &&
\end{calculs}
\explique{On retrouve $\res{\vect{AC}}\coord{4-2}{-2-1}$, soit $\vect{AC}\coord{2}{-3}$.}{relation de
Chasles : $\vect{AB}+\vect{BC}=\vect{AC}$}
\cas{c)}
\explique{On pose $\vec{w}\coord{x}{y}$.}{les coordonnées cherchées sont les inconnues}
\begin{calculs}
\vec{u}+\vec{w} & \coord{3+x}{-2+y} &\qquad& \rem{abscisses ensemble, ordonnées ensemble}\\
3+x &= -1 && \rem{même abscisse que $\vec{v}$}\\
x &= -4 && \rem{on retranche $3$}\\
-2+y &= 5 && \rem{même ordonnée que $\vec{v}$}\\
y &= 7 && \rem{on ajoute $2$}
\end{calculs}
\explique{Le vecteur $\vec{w}$ a pour coordonnées $\res{\vec{w}\coord{-4}{7}}$.}{on vérifie :
c'est $\vec{v}-\vec{u}$}
\end{correction}

\etape{À vous de jouer}
On considère les vecteurs $\vec{u}\coord{-4}{1}$ et $\vec{v}\coord{2}{3}$, et les points
$E(1\,;-2)$, $F(3\,;2)$ et $G(-2\,;0)$.
\begin{questions}[label=\alph*)]
  \item Calculer les coordonnées des vecteurs $\vec{u}+\vec{v}$ et $\vec{u}-\vec{v}$.
  \item Calculer les coordonnées du vecteur $\vect{EF}+\vect{FG}$. Quel vecteur
        retrouve-t-on ?
  \item Calculer les coordonnées du vecteur $\vec{w}$ tel que $\vec{u}+\vec{w}=\vec{v}$.
\end{questions}
\reponse{a)~$\vec{u}+\vec{v}\coord{-2}{4}$ et $\vec{u}-\vec{v}\coord{-6}{-2}$ ;\;
b)~$\coord{-3}{2}$, c'est $\vect{EG}$ ;\; c)~$\vec{w}\coord{6}{2}$.}
\end{methodebox}

\afaire{exercices 1 à 4}

\partiecours{Parallélogramme et milieu d'un segment}

\begin{proprietebox}[Propriété 3 --- Égalité de vecteurs et parallélogramme]
\begin{minipage}[t]{0.58\linewidth}\raggedright
\begin{itemize}
  \item Deux vecteurs sont \textbf{égaux} quand ils ont les mêmes coordonnées.
  \item $ABCD$ est un \textbf{parallélogramme} quand $\vect{AB}=\vect{DC}$.
  \item $ABCD$ est un parallélogramme quand ses diagonales $[AC]$ et $[BD]$ ont le
        \textbf{même milieu}.
\end{itemize}
\end{minipage}\hfill
\begin{minipage}[t]{0.38\linewidth}
\vspace{0pt}\centering
\begin{tikzpicture}[scale=0.55]
  \draw[fig] (0,0) -- (4,1) -- (5,3) -- (1,2) -- cycle;
  \draw[aide] (0,0) -- (5,3) (4,1) -- (1,2);
  \draw[vecteur] (0,0) -- node[midway,below right,font=\footnotesize]{$\vect{AB}$} (4,1);
  \draw[vecteurB] (1,2) -- node[midway,above left,font=\footnotesize]{$\vect{DC}$} (5,3);
  \pt{0,0}{below left}{A} \pt{4,1}{below right}{B} \pt{5,3}{above right}{C}
  \pt{1,2}{above left}{D}
\end{tikzpicture}
\end{minipage}
\end{proprietebox}

\begin{proprietebox}[Propriété 4 --- Coordonnées du milieu]
Le milieu $I$ du segment $[AB]$ a pour coordonnées les \textbf{moyennes} des
coordonnées de $A$ et de $B$ :
\begin{center}
  $x_{I}=\dfrac{x_{A}+x_{B}}{2}$ \qquad $y_{I}=\dfrac{y_{A}+y_{B}}{2}$
\end{center}
\demo{$I$ est le milieu de $[AB]$, donc $\vect{AI}=\vect{IB}$ : les abscisses de ces deux
vecteurs sont égales.}
\begin{calculs}
x_{I}-x_{A} &= x_{B}-x_{I} &\qquad& \rem{arrivée moins départ, des deux côtés}\\
2x_{I} &= x_{A}+x_{B} && \rem{on ajoute $x_{I}+x_{A}$}\\
x_{I} &= \dfrac{x_{A}+x_{B}}{2} && \rem{de même pour les ordonnées}
\end{calculs}
\end{proprietebox}

\begin{methodebox}[Méthode 3 --- Calculer les coordonnées d'un point défini par une égalité de vecteurs]
L'égalité de deux vecteurs donne deux équations : une pour les abscisses, une pour
les ordonnées.
\etapeexemples
On considère les points $A(-1\,;2)$, $B(3\,;4)$, $C(4\,;1)$ et le vecteur
$\vec{u}\coord{2}{-5}$. Calculer les coordonnées :
\begin{questions}[label=\alph*)]
  \item du point $E$ tel que $\vect{AE}=\vec{u}$ ;
  \item du point $D$ tel que $ABCD$ soit un parallélogramme ;
  \item du point $F$ tel que $\vect{FA}=\vect{BC}$.
\end{questions}
\begin{correction}
\cas{a)}
\begin{calculs}
\vect{AE} & \coord{x_{E}-(-1)}{y_{E}-2} &\qquad& \rem{arrivée $E$ (inconnue) moins départ $A$}\\
x_{E}+1 &= 2 && \rem{même abscisse que $\vec{u}$}\\
x_{E} &= 1 && \rem{on retranche $1$}\\
y_{E}-2 &= -5 && \rem{même ordonnée que $\vec{u}$}\\
y_{E} &= -3 && \rem{on ajoute $2$}
\end{calculs}
\explique{Le point $E$ a pour coordonnées $\res{(1\,;-3)}$.}{on conclut}
\cas{b)}
\explique{$ABCD$ est un parallélogramme quand $\vect{AB}=\vect{DC}$.}{propriété 3 ;
l'inconnue est $D$}
\begin{calculs}
\vect{AB} & \coord{3-(-1)}{4-2} &\qquad& \rem{arrivée moins départ}\\
\vect{AB} & \coord{4}{2} &&\\
\vect{DC} & \coord{4-x_{D}}{1-y_{D}} && \rem{arrivée $C$ moins départ $D$}\\
4-x_{D} &= 4 && \rem{les abscisses sont égales}\\
x_{D} &= 0 &&\\
1-y_{D} &= 2 && \rem{les ordonnées sont égales}\\
-y_{D} &= 1 && \rem{on retranche $1$}\\
y_{D} &= -1 && \rem{on change les signes}
\end{calculs}
\explique{Le point $D$ a pour coordonnées $\res{(0\,;-1)}$.}{on conclut}
\cas{c)}
\begin{calculs}
\vect{BC} & \coord{4-3}{1-4} &\qquad& \rem{le vecteur connu}\\
\vect{BC} & \coord{1}{-3} &&\\
\vect{FA} & \coord{-1-x_{F}}{2-y_{F}} && \rem{l'inconnue $F$ est au départ}\\
-1-x_{F} &= 1 && \rem{les abscisses sont égales}\\
-x_{F} &= 2 && \rem{on ajoute $1$}\\
x_{F} &= -2 &&\\
2-y_{F} &= -3 && \rem{les ordonnées sont égales}\\
-y_{F} &= -5 && \rem{on retranche $2$}\\
y_{F} &= 5 &&
\end{calculs}
\explique{Le point $F$ a pour coordonnées $\res{(-2\,;5)}$.}{on conclut}
\end{correction}

\etape{À vous de jouer}
On considère les points $K(2\,;-1)$, $L(-1\,;3)$, $M(1\,;4)$ et le vecteur
$\vec{w}\coord{-3}{2}$. Calculer les coordonnées :
\begin{questions}[label=\alph*)]
  \item du point $N$ tel que $\vect{KN}=\vec{w}$ ;
  \item du point $R$ tel que $KLMR$ soit un parallélogramme ;
  \item du point $S$ tel que $\vect{SK}=\vect{LM}$.
\end{questions}
\reponse{a)~$N(-1\,;1)$ ;\; b)~$R(4\,;0)$ ;\; c)~$S(0\,;-2)$.}
\end{methodebox}

\begin{methodebox}[Méthode 4 --- Utiliser les coordonnées du milieu]
\etapeexemples
\begin{questions}[label=\alph*)]
  \item On considère les points $A(-3\,;4)$ et $B(5\,;-2)$. Calculer les coordonnées du
        milieu $I$ de $[AB]$.
  \item Le point $K(2\,;-1)$ est le milieu du segment $[AC]$. Calculer les coordonnées
        du point $C$.
  \item On considère les points $E(-2\,;1)$, $F(3\,;2)$, $G(4\,;-1)$ et $H(-1\,;-2)$.
        Démontrer que $EFGH$ est un parallélogramme.
\end{questions}
\begin{correction}
\cas{a)}
\begin{calculs}
x_{I} &= \dfrac{x_{A}+x_{B}}{2} &\qquad& \rem{on recopie la formule}\\
x_{I} &= \dfrac{-3+5}{2} && \rem{on remplace}\\
x_{I} &= 1 &&\\
y_{I} &= \dfrac{4+(-2)}{2} && \rem{de même pour l'ordonnée}\\
y_{I} &= 1 &&
\end{calculs}
\explique{Le milieu de $[AB]$ est $\res{I(1\,;1)}$.}{on conclut}
\cas{b)}
\explique{L'inconnue est $C$ : on écrit les coordonnées du milieu $K$ avec $x_{C}$ et
$y_{C}$.}{$K$ est connu, $C$ ne l'est pas}
\begin{calculs}
\dfrac{-3+x_{C}}{2} &= 2 &\qquad& \rem{l'abscisse de $K$}\\
-3+x_{C} &= 4 && \rem{on multiplie par $2$}\\
x_{C} &= 7 && \rem{on ajoute $3$}\\
\dfrac{4+y_{C}}{2} &= -1 && \rem{l'ordonnée de $K$}\\
4+y_{C} &= -2 && \rem{on multiplie par $2$}\\
y_{C} &= -6 && \rem{on retranche $4$}
\end{calculs}
\explique{Le point $C$ a pour coordonnées $\res{(7\,;-6)}$.}{on vérifie : $\frac{-3+7}{2}=2$
et $\frac{4-6}{2}=-1$}
\cas{c)}
\explique{On calcule le milieu $J$ de la diagonale $[EG]$.}{propriété 3 : les deux
diagonales}
\begin{calculs}
x_{J} &= \dfrac{-2+4}{2} &\qquad& \rem{moyenne des abscisses}\\
x_{J} &= 1 &&\\
y_{J} &= \dfrac{1+(-1)}{2} && \rem{moyenne des ordonnées}\\
y_{J} &= 0 &&
\end{calculs}
\explique{On calcule le milieu $J'$ de l'autre diagonale $[FH]$.}{même méthode}
\begin{calculs}
x_{J'} &= \dfrac{3+(-1)}{2} &\qquad& \rem{moyenne des abscisses}\\
x_{J'} &= 1 &&\\
y_{J'} &= \dfrac{2+(-2)}{2} && \rem{moyenne des ordonnées}\\
y_{J'} &= 0 &&
\end{calculs}
\explique{Les diagonales $[EG]$ et $[FH]$ ont le même milieu $(1\,;0)$ : $EFGH$ \res{est un
parallélogramme}.}{propriété 3}
\end{correction}

\etape{À vous de jouer}
\begin{questions}[label=\alph*)]
  \item On considère les points $R(-5\,;2)$ et $S(3\,;6)$. Calculer les coordonnées du
        milieu $J$ de $[RS]$.
  \item Le point $M(0\,;-1)$ est le milieu du segment $[RT]$. Calculer les coordonnées
        du point $T$.
  \item On considère les points $K(-3\,;0)$, $L(1\,;3)$, $U(5\,;2)$ et $V(1\,;-1)$.
        Démontrer que $KLUV$ est un parallélogramme.
\end{questions}
\reponse{a)~$J(-1\,;4)$ ;\; b)~$T(5\,;-4)$ ;\; c)~$[KU]$ et $[LV]$ ont le même milieu
$(1\,;1)$.}
\end{methodebox}

\afaire{exercices 5 à 8}

\partiecours{Produit d'un vecteur par un nombre réel, colinéarité}

\begin{definitionbox}[Définition 2 --- Produit d'un vecteur par un nombre réel]
\begin{minipage}[t]{0.58\linewidth}\raggedright
Pour un vecteur non nul $\vec{u}$ et un nombre réel $k$ non nul, le vecteur $k\vec{u}$ a :
\begin{itemize}
  \item la même \textbf{direction} que $\vec{u}$ ;
  \item le même \textbf{sens} si $k>0$, le sens contraire si $k<0$ ;
  \item une \textbf{longueur} multipliée par $k$ si $k>0$, par $-k$ si $k<0$.
\end{itemize}
Si $k=0$ ou $\vec{u}=\vec{0}$, alors $k\vec{u}=\vec{0}$.
\end{minipage}\hfill
\begin{minipage}[t]{0.38\linewidth}
\vspace{0pt}\centering
\begin{tikzpicture}[x=0.55cm,y=0.55cm]
  \quadrillage{-0.4}{-0.4}{6.4}{3.4}
  \draw[vecteur] (0,1) -- node[midway,above left,font=\footnotesize]{$\vec{u}$} (2,2);
  \draw[vecteurB] (2,0) -- node[midway,below right,font=\footnotesize]{$2\vec{u}$} (6,2);
  \draw[vecteurC] (3,3) -- node[midway,above left,font=\footnotesize]{$-\vec{u}$} (1,2);
\end{tikzpicture}
\end{minipage}

\medskip
En coordonnées : pour $\vec{u}\coord{x}{y}$, on a $k\vec{u}\coord{kx}{ky}$. Par exemple, pour
$\vec{u}\coord{2}{-1}$ et $\vec{v}\coord{1}{4}$ :
\begin{center}
  $2\vec{u}\coord{4}{-2}$ \qquad $3\vec{v}\coord{3}{12}$ \qquad
  $2\vec{u}-3\vec{v}\coord{4-3}{-2-12}$, soit $2\vec{u}-3\vec{v}\coord{1}{-14}$.
\end{center}
\end{definitionbox}

\begin{definitionbox}[Définition 3 --- Vecteurs colinéaires, déterminant]
Deux vecteurs $\vec{u}$ et $\vec{v}$ sont \textbf{colinéaires} quand l'un est le produit de
l'autre par un nombre réel : $\vec{v}=k\vec{u}$ (ou $\vec{u}=k\vec{v}$). Le vecteur nul est
colinéaire à tout vecteur.

\smallskip
Le \textbf{déterminant} de $\vec{u}\coord{x}{y}$ et $\vec{v}\coord{x'}{y'}$ est le nombre
\[
  \det(\vec{u}\,;\vec{v})=\begin{vmatrix} x & x'\\ y & y'\end{vmatrix}=xy'-yx'
  \qquad\text{(produit en croix des coordonnées).}
\]
\end{definitionbox}

\begin{proprietebox}[Théorème 1 --- Critère de colinéarité]
$\vec{u}$ et $\vec{v}$ sont colinéaires quand $\det(\vec{u}\,;\vec{v})=0$ ; ils ne le
sont pas quand $\det(\vec{u}\,;\vec{v})\neq 0$.

\smallskip
\demo{d'un sens (la réciproque est admise) : si $\vec{v}=k\vec{u}$, alors $x'=kx$ et
$y'=ky$.}
\begin{calculs}
\det(\vec{u}\,;\vec{v}) &= x\times ky-y\times kx &\qquad& \rem{$xy'-yx'$}\\
\det(\vec{u}\,;\vec{v}) &= 0 && \rem{les deux produits sont égaux}
\end{calculs}
\end{proprietebox}

\begin{methodebox}[Méthode 5 --- Démontrer que deux vecteurs sont colinéaires, ou non]
On cherche d'abord un multiple évident ($\vec{v}=k\vec{u}$) ; s'il ne se voit pas, on
calcule le déterminant.
\etapeexemples
\begin{questions}[label=\alph*)]
  \item Les vecteurs $\vec{u}\coord{2}{-3}$ et $\vec{v}\coord{-6}{9}$ sont-ils colinéaires ?
  \item Les vecteurs $\vec{u}\coord{3}{5}$ et $\vec{w}\coord{5}{8}$ sont-ils colinéaires ?
  \item Les vecteurs $\vec{a}\coord{\sqrt{2}}{1}$ et $\vec{b}\coord{2}{\sqrt{2}}$ sont-ils
        colinéaires ?
  \item Calculer les nombres réels $m$ pour lesquels les vecteurs $\vec{r}\coord{m}{2}$
        et $\vec{s}\coord{8}{m}$ sont colinéaires.
\end{questions}
\begin{correction}
\cas{a)}
\explique{$-6=(-3)\times 2$ et $9=(-3)\times(-3)$, donc $\vec{v}=-3\vec{u}$.}{même
multiplicateur pour l'abscisse et pour l'ordonnée}
\explique{Les vecteurs $\vec{u}$ et $\vec{v}$ \res{sont colinéaires}.}{de sens contraires :
$-3<0$}
\cas{b)}
\begin{calculs}
\det(\vec{u}\,;\vec{w}) &= xy'-yx' &\qquad& \rem{pas de multiple évident : on recopie la formule}\\
\det(\vec{u}\,;\vec{w}) &= 3\times 8-5\times 5 && \rem{$x=3$ et $y=5$, puis $x'=5$ et $y'=8$}\\
\det(\vec{u}\,;\vec{w}) &= -1 && \rem{$24-25$}
\end{calculs}
\explique{Le déterminant n'est pas nul : $\vec{u}$ et $\vec{w}$ \res{ne sont pas
colinéaires}.}{$-1\neq 0$}
\cas{c)}
\begin{calculs}
\det(\vec{a}\,;\vec{b}) &= \sqrt{2}\times\sqrt{2}-1\times 2 &\qquad& \rem{$xy'-yx'$}\\
\det(\vec{a}\,;\vec{b}) &= 2-2 && \rem{$\sqrt{2}\times\sqrt{2}=2$}\\
\det(\vec{a}\,;\vec{b}) &= 0 &&
\end{calculs}
\explique{Le déterminant est nul : $\vec{a}$ et $\vec{b}$ \res{sont colinéaires}.}{en effet,
$\vec{b}=\sqrt{2}\,\vec{a}$}
\cas{d)}
\begin{calculs}
\det(\vec{r}\,;\vec{s}) &= m\times m-2\times 8 &\qquad& \rem{$xy'-yx'$}\\
\det(\vec{r}\,;\vec{s}) &= m^{2}-16 &&
\end{calculs}
\explique{Les vecteurs sont colinéaires quand ce déterminant est nul.}{on écrit l'équation}
\begin{calculs}
m^{2}-16 &= 0 &\qquad& \rem{l'équation à résoudre}\\
(m-4)(m+4) &= 0 && \rem{$a^{2}-b^{2}=(a-b)(a+b)$}
\end{calculs}
\explique{Un produit est nul quand l'un de ses facteurs est nul : $\res{m=4}$ ou
$\res{m=-4}$.}{on vérifie : pour $m=4$, $\vec{s}=2\vec{r}$}
\end{correction}

\etape{À vous de jouer}
\begin{questions}[label=\alph*)]
  \item Les vecteurs $\vec{u}\coord{-1}{4}$ et $\vec{v}\coord{3}{-12}$ sont-ils colinéaires ?
  \item Les vecteurs $\vec{u}\coord{4}{7}$ et $\vec{w}\coord{5}{9}$ sont-ils colinéaires ?
  \item Les vecteurs $\vec{a}\coord{\sqrt{3}}{3}$ et $\vec{b}\coord{1}{\sqrt{3}}$ sont-ils
        colinéaires ?
  \item Calculer les nombres réels $m$ pour lesquels les vecteurs $\vec{r}\coord{m}{3}$
        et $\vec{s}\coord{12}{m}$ sont colinéaires.
\end{questions}
\reponse{a)~oui, $\vec{v}=-3\vec{u}$ ;\; b)~non, le déterminant vaut $1$ ;\; c)~oui, le
déterminant est nul ;\; d)~$m=6$ ou $m=-6$.}
\end{methodebox}

\begin{proprietebox}[Propriété 5 --- Parallélisme et alignement]
\begin{itemize}
  \item Les droites $(AB)$ et $(CD)$ sont \textbf{parallèles} quand les vecteurs
        $\vect{AB}$ et $\vect{CD}$ sont colinéaires.
  \item Les points $A$, $B$ et $C$ sont \textbf{alignés} quand les vecteurs
        $\vect{AB}$ et $\vect{AC}$ sont colinéaires.
\end{itemize}
\end{proprietebox}

\begin{methodebox}[Méthode 6 --- Démontrer un parallélisme ou un alignement]
\etapeexemples
On considère les points $A(-1\,;-2)$, $B(1\,;2)$, $C(3\,;0)$, $D(4\,;2)$ et $E(3\,;6)$.
\begin{questions}[label=\alph*)]
  \item Les droites $(AB)$ et $(CD)$ sont-elles parallèles ?
  \item Les points $A$, $B$ et $E$ sont-ils alignés ?
  \item Les points $A$, $B$ et $C$ sont-ils alignés ?
  \item Calculer l'ordonnée $y$ du point $M(4\,;y)$ de la droite $(AB)$.
\end{questions}
\begin{correction}
\cas{a)}
\begin{calculs}
\vect{AB} & \coord{1-(-1)}{2-(-2)} &\qquad& \rem{on calcule les deux vecteurs}\\
\vect{AB} & \coord{2}{4} &&\\
\vect{CD} & \coord{4-3}{2-0} &&\\
\vect{CD} & \coord{1}{2} &&\\
\det(\vect{AB}\,;\vect{CD}) &= 2\times 2-4\times 1 && \rem{puis leur déterminant}\\
\det(\vect{AB}\,;\vect{CD}) &= 0 &&
\end{calculs}
\explique{$\vect{AB}$ et $\vect{CD}$ sont colinéaires : les droites $(AB)$ et $(CD)$
\res{sont parallèles}.}{propriété 5}
\cas{b)}
\begin{calculs}
\vect{AE} & \coord{3-(-1)}{6-(-2)} &\qquad& \rem{$\vect{AB}\coord{2}{4}$ est déjà calculé}\\
\vect{AE} & \coord{4}{8} &&\\
\det(\vect{AB}\,;\vect{AE}) &= 2\times 8-4\times 4 &&\\
\det(\vect{AB}\,;\vect{AE}) &= 0 &&
\end{calculs}
\explique{$\vect{AB}$ et $\vect{AE}$ sont colinéaires : les points $A$, $B$ et $E$ \res{sont
alignés}.}{propriété 5}
\cas{c)}
\begin{calculs}
\vect{AC} & \coord{3-(-1)}{0-(-2)} &\qquad& \rem{arrivée moins départ}\\
\vect{AC} & \coord{4}{2} &&\\
\det(\vect{AB}\,;\vect{AC}) &= 2\times 2-4\times 4 &&\\
\det(\vect{AB}\,;\vect{AC}) &= -12 &&
\end{calculs}
\explique{Le déterminant n'est pas nul : $A$, $B$ et $C$ \res{ne sont pas alignés}.}{$-12\neq 0$}
\cas{d)}
\explique{$M$ est sur la droite $(AB)$ quand $\vect{AB}$ et $\vect{AM}$ sont
colinéaires.}{l'inconnue est $y$}
\begin{calculs}
\vect{AM} & \coord{4-(-1)}{y-(-2)} &\qquad& \rem{arrivée moins départ}\\
\vect{AM} & \coord{5}{y+2} &&\\
\det(\vect{AB}\,;\vect{AM}) &= 2(y+2)-4\times 5 && \rem{$xy'-yx'$}\\
\det(\vect{AB}\,;\vect{AM}) &= 2y-16 && \rem{$2y+4-20$}\\
2y-16 &= 0 && \rem{le déterminant doit être nul}\\
y &= 8 && \rem{on ajoute $16$, puis on divise par $2$}
\end{calculs}
\explique{Le point $M$ a pour coordonnées $\res{(4\,;8)}$.}{on conclut}
\end{correction}

\etape{À vous de jouer}
On considère les points $F(-2\,;1)$, $G(0\,;2)$, $H(1\,;-3)$, $K(5\,;-1)$ et $L(6\,;5)$.
\begin{questions}[label=\alph*)]
  \item Les droites $(FG)$ et $(HK)$ sont-elles parallèles ?
  \item Les points $F$, $G$ et $L$ sont-ils alignés ?
  \item Les points $F$, $G$ et $H$ sont-ils alignés ?
  \item Calculer l'ordonnée $y$ du point $N(8\,;y)$ de la droite $(FG)$.
\end{questions}
\reponse{a)~oui ;\; b)~oui ;\; c)~non, le déterminant vaut $-11$ ;\; d)~$y=6$.}
\end{methodebox}

\afaire{exercices 9 à 12}

\partiecours{Norme d'un vecteur, équation d'un cercle}

À partir d'ici, le repère $(O\,;\vec{\imath},\vec{\jmath})$ est \textbf{orthonormé} :
$\vec{\imath}$ et $\vec{\jmath}$ ont des directions perpendiculaires et la même
longueur, prise pour unité.

\begin{proprietebox}[Propriété 6 --- Norme et distance]
\begin{minipage}[t]{0.58\linewidth}\raggedright
La \textbf{norme} d'un vecteur $\vec{u}$, notée $\norme{\vec{u}}$, est sa longueur ; pour
deux points, $\norme{\vect{AB}}=AB$. Pour $\vec{u}\coord{x}{y}$ :
\[
  \norme{\vec{u}}=\sqrt{x^{2}+y^{2}}.
\]
Pour deux points $A$ et $B$ :
$AB=\sqrt{(x_{B}-x_{A})^{2}+(y_{B}-y_{A})^{2}}$.
\end{minipage}\hfill
\begin{minipage}[t]{0.38\linewidth}
\vspace{0pt}\centering
\begin{tikzpicture}[x=0.6cm,y=0.6cm]
  \quadrillage{-0.4}{-0.4}{4.4}{3.4}
  \draw[aide] (0,0) -- (4,0) -- (4,3);
  \draw[black!60] (3.7,0) -- (3.7,0.3) -- (4,0.3);
  \draw[vecteur] (0,0) -- node[midway,above left,font=\footnotesize]{$\vec{u}$} (4,3);
  \node[below,font=\footnotesize] at (2,0) {$x$};
  \node[right,font=\footnotesize] at (4,1.5) {$y$};
\end{tikzpicture}
\end{minipage}

\smallskip
\demo{le triangle de la figure est rectangle (le repère est orthonormé) ; ses côtés de
l'angle droit mesurent $x$ et $y$, au signe près, et le théorème de Pythagore donne
$\norme{\vec{u}}^{2}=x^{2}+y^{2}$.}
\end{proprietebox}

\begin{methodebox}[Méthode 7 --- Calculer une norme]
\etapeexemples
\begin{questions}[label=\alph*)]
  \item Calculer la norme du vecteur $\vec{u}\coord{-5}{12}$.
  \item On considère les points $A(-2\,;1)$ et $B(2\,;3)$. Calculer la distance $AB$
        et l'écrire sous la forme $a\sqrt{b}$, avec deux nombres entiers $a$ et $b$.
  \item On considère les points $C(-1\,;0)$, $D(3\,;2)$ et $E(1\,;6)$. Quelle est la
        nature du triangle $CDE$ ?
\end{questions}
\begin{correction}
\cas{a)}
\begin{calculs}
\norme{\vec{u}} &= \sqrt{x^{2}+y^{2}} &\qquad& \rem{on recopie la formule}\\
\norme{\vec{u}} &= \sqrt{(-5)^{2}+12^{2}} && \rem{parenthèses autour de $-5$}\\
\norme{\vec{u}} &= \sqrt{169} && \rem{$25+144$}\\
\norme{\vec{u}} &= \res{13} &&
\end{calculs}
\cas{b)}
\begin{calculs}
\vect{AB} & \coord{4}{2} &\qquad& \rem{$2-(-2)$ et $3-1$}\\
AB &= \sqrt{4^{2}+2^{2}} && \rem{la norme de $\vect{AB}$}\\
AB &= \sqrt{20} &&\\
AB &= \sqrt{4}\times\sqrt{5} && \rem{$20=4\times 5$, et $4$ est un carré}\\
AB &= \res{2\sqrt{5}} &&
\end{calculs}
\cas{c)}
\explique{On calcule les carrés des trois côtés : pas de racine carrée.}{Pythagore
s'écrit avec des carrés}
\begin{calculs}
\vect{CD} & \coord{4}{2} &\qquad& \rem{$3-(-1)$ et $2-0$}\\
CD^{2} &= 20 && \rem{$4^{2}+2^{2}$}\\
\vect{CE} & \coord{2}{6} && \rem{$1-(-1)$ et $6-0$}\\
CE^{2} &= 40 && \rem{$2^{2}+6^{2}$}\\
\vect{DE} & \coord{-2}{4} && \rem{$1-3$ et $6-2$}\\
DE^{2} &= 20 && \rem{$(-2)^{2}+4^{2}$}
\end{calculs}
\explique{$CD=DE$, et $CD^{2}+DE^{2}=40=CE^{2}$ : le triangle $CDE$ \res{est rectangle et
isocèle en $D$}.}{réciproque du théorème de Pythagore}
\end{correction}

\etape{À vous de jouer}
\begin{questions}[label=\alph*)]
  \item Calculer la norme du vecteur $\vec{v}\coord{8}{-15}$.
  \item On considère les points $E(-3\,;2)$ et $F(3\,;-1)$. Calculer la distance $EF$
        et l'écrire sous la forme $a\sqrt{b}$, avec deux nombres entiers $a$ et $b$.
  \item On considère les points $K(-2\,;1)$, $L(1\,;2)$ et $M(0\,;5)$. Quelle est la
        nature du triangle $KLM$ ?
\end{questions}
\reponse{a)~$17$ ;\; b)~$EF=3\sqrt{5}$ ;\; c)~rectangle et isocèle en $L$.}
\end{methodebox}

\begin{proprietebox}[Théorème 2 --- Équation d'un cercle]
\begin{minipage}[t]{0.58\linewidth}\raggedright
Le cercle de centre $\Omega(a\,;b)$ et de rayon $r$ est l'ensemble des points $M(x\,;y)$
tels que
\[
  (x-a)^{2}+(y-b)^{2}=r^{2}.
\]
On dit que c'est une \textbf{équation} du cercle.
\end{minipage}\hfill
\begin{minipage}[t]{0.38\linewidth}
\vspace{0pt}\centering
\begin{tikzpicture}[scale=0.55]
  \draw[coulCourbe,line width=1pt] (0,0) circle (2);
  \draw[aide] (0,0) -- (1.6,1.2);
  \node[font=\footnotesize,above left] at (0.8,0.6) {$r$};
  \pt{0,0}{below}{\Omega} \pt{1.6,1.2}{above right}{M}
\end{tikzpicture}
\end{minipage}

\smallskip
\demo{une distance est positive, donc deux distances sont égales quand leurs carrés le
sont.}
\begin{calculs}
\Omega M &= r &\qquad& \rem{$M$ est sur le cercle}\\
\Omega M^{2} &= r^{2} && \rem{on passe aux carrés}\\
(x-a)^{2}+(y-b)^{2} &= r^{2} && \rem{propriété 6, avec $\vect{\Omega M}\coord{x-a}{y-b}$}
\end{calculs}
\smallskip
En développant, on obtient une autre forme : $(x-1)^{2}+(y+3)^{2}=4$ devient
$x^{2}+y^{2}-2x+6y+6=0$. Pour reconnaître un cercle sous cette forme, on fait le chemin
inverse en complétant les carrés (forme canonique, chapitre 3) : voir la fiche
d'exercices.
\end{proprietebox}

\begin{methodebox}[Méthode 8 --- Trouver l'équation cartésienne d'un cercle]
\etapeexemples
\begin{questions}[label=\alph*)]
  \item Lire le centre et le rayon du cercle d'équation $(x+3)^{2}+(y-1)^{2}=7$.
  \item Écrire une équation du cercle $\mathcal{C}_{1}$ de centre $\Omega(2\,;-3)$ et de rayon $4$.
  \item Écrire une équation du cercle $\mathcal{C}_{2}$ de centre $A(-1\,;2)$ qui passe par
        $B(3\,;5)$.
  \item Les points $C(2\,;-2)$ et $D(4\,;1)$ sont-ils sur le cercle $\mathcal{C}_{2}$ ?
  \item Écrire une équation du cercle $\mathcal{C}_{3}$ de diamètre $[EF]$, avec $E(-3\,;1)$ et
        $F(5\,;-5)$.
\end{questions}
\begin{correction}
\cas{a)}
\explique{$(x+3)^{2}=\big(x-(-3)\big)^{2}$ et $(y-1)^{2}$ : le centre est
$\res{(-3\,;1)}$.}{on lit $a$ et $b$ dans $(x-a)^{2}+(y-b)^{2}$ : attention aux signes}
\explique{$r^{2}=7$ et $r>0$, donc le rayon vaut $r=\res{\sqrt{7}}$.}{le membre de droite
est le carré du rayon}
\cas{b)}
\begin{calculs}
(x-a)^{2}+(y-b)^{2} &= r^{2} &\qquad& \rem{on recopie l'équation}\\
(x-2)^{2}+\big(y-(-3)\big)^{2} &= 4^{2} && \rem{$a=2$, $b=-3$ et $r=4$}\\
\res{(x-2)^{2}+(y+3)^{2}} &= \res{16} &&
\end{calculs}
\cas{c)}
\explique{Le rayon est la distance $AB$ : on calcule $r^{2}=AB^{2}$.}{pas besoin de la
racine carrée}
\begin{calculs}
\vect{AB} & \coord{4}{3} &\qquad& \rem{$3-(-1)$ et $5-2$}\\
AB^{2} &= 25 && \rem{$4^{2}+3^{2}$}
\end{calculs}
\explique{Une équation de $\mathcal{C}_{2}$ est
$\res{(x+1)^{2}+(y-2)^{2}=25}$.}{$a=-1$ et $b=2$}
\cas{d)}
\begin{calculs}
(x_{C}+1)^{2}+(y_{C}-2)^{2} &= (2+1)^{2}+(-2-2)^{2} &\qquad& \rem{on remplace par les coordonnées de $C$}\\
(x_{C}+1)^{2}+(y_{C}-2)^{2} &= 25 && \rem{$9+16$}\\
(x_{D}+1)^{2}+(y_{D}-2)^{2} &= (4+1)^{2}+(1-2)^{2} && \rem{de même pour $D$}\\
(x_{D}+1)^{2}+(y_{D}-2)^{2} &= 26 && \rem{$25+1$}
\end{calculs}
\explique{Les coordonnées de $C$ vérifient l'équation : $C$ \res{est sur $\mathcal{C}_{2}$} ; $D$
\res{n'y est pas}.}{$26\neq 25$}
\cas{e)}
\explique{Le centre est le milieu $\Omega_{3}$ de $[EF]$ et le rayon vaut
$\frac{EF}{2}$.}{un diamètre passe par le centre}
\begin{calculs}
x_{\Omega_{3}} &= \dfrac{-3+5}{2} &\qquad& \rem{moyenne des abscisses}\\
x_{\Omega_{3}} &= 1 &&\\
y_{\Omega_{3}} &= \dfrac{1+(-5)}{2} && \rem{moyenne des ordonnées}\\
y_{\Omega_{3}} &= -2 &&\\
\vect{EF} & \coord{8}{-6} && \rem{$5-(-3)$ et $-5-1$}\\
EF &= 10 && \rem{$\sqrt{64+36}$}\\
r &= 5 && \rem{la moitié de $EF$}
\end{calculs}
\explique{Une équation de $\mathcal{C}_{3}$ est
$\res{(x-1)^{2}+(y+2)^{2}=25}$.}{$r^{2}=25$}
\end{correction}

\etape{À vous de jouer}
\begin{questions}[label=\alph*)]
  \item Lire le centre et le rayon du cercle d'équation $(x-5)^{2}+(y+2)^{2}=11$.
  \item Écrire une équation du cercle $\mathcal{C}_{1}$ de centre $\Omega(-4\,;1)$ et de rayon $3$.
  \item Écrire une équation du cercle $\mathcal{C}_{2}$ de centre $K(3\,;-1)$ qui passe par
        $L(-1\,;2)$.
  \item Les points $M(6\,;3)$ et $N(0\,;2)$ sont-ils sur le cercle $\mathcal{C}_{2}$ ?
  \item Écrire une équation du cercle $\mathcal{C}_{3}$ de diamètre $[GH]$, avec $G(-2\,;-3)$ et
        $H(4\,;5)$.
\end{questions}
\reponse{a)~centre $(5\,;-2)$ et rayon $\sqrt{11}$ ;\; b)~$(x+4)^{2}+(y-1)^{2}=9$ ;\;
c)~$(x-3)^{2}+(y+1)^{2}=25$ ;\; d)~$M$ oui, $N$ non ;\; e)~$(x-1)^{2}+(y-1)^{2}=25$.}
\end{methodebox}

\afaire{exercices 13 à 16}

\partiecours{Définition du produit scalaire}

\begin{definitionbox}[Définition 4 --- Produit scalaire]
\begin{minipage}[t]{0.58\linewidth}\raggedright
On représente deux vecteurs non nuls avec la même origine : $\vec{u}=\vect{AB}$ et
$\vec{v}=\vect{AC}$. Leur \textbf{angle} est l'angle géométrique $\widehat{BAC}$, entre
$0^{\circ}$ et $180^{\circ}$. Le \textbf{produit scalaire} de $\vec{u}$ et $\vec{v}$ est le
\textbf{nombre}
\[
  \vec{u}\cdot\vec{v}=\norme{\vec{u}}\times\norme{\vec{v}}\times\cos\big(\widehat{BAC}\big),
\]
c'est-à-dire $\vect{AB}\cdot\vect{AC}=AB\times AC\times\cos\big(\widehat{BAC}\big)$. Il se lit
« $\vec{u}$ scalaire $\vec{v}$ ». Si l'un des deux vecteurs est nul,
$\vec{u}\cdot\vec{v}=0$.
\end{minipage}\hfill
\begin{minipage}[t]{0.38\linewidth}
\vspace{0pt}\centering
\begin{tikzpicture}[scale=0.8]
  \draw (0.7,0) arc (0:38.66:0.7);
  \draw[vecteur] (0,0) -- node[midway,below,font=\footnotesize]{$\vec{u}$} (3.5,0);
  \draw[vecteurB] (0,0) -- node[midway,above left,font=\footnotesize]{$\vec{v}$} (2.5,2);
  \pt{0,0}{left}{A} \pt{3.5,0}{below right}{B} \pt{2.5,2}{above right}{C}
\end{tikzpicture}
\end{minipage}

\smallskip
\textbf{Carré scalaire} : $\vec{u}\cdot\vec{u}=\norme{\vec{u}}^{2}$ (l'angle est nul) ; on
note $\vec{u}^{2}=\vec{u}\cdot\vec{u}$. Par exemple, $\vect{AB}\cdot\vect{AB}=AB^{2}$.
\end{definitionbox}

\begin{remarquebox}[Cosinus d'un angle (rappel du chapitre 8)]
Un angle géométrique se mesure ici en \textbf{degrés} ; $180^{\circ}$ correspond à
$\pi$ radians. Son cosinus se lit sur le cercle trigonométrique (abscisse du point
associé) ; à la calculatrice, en mode \textbf{degré}. Valeurs remarquables :
\begin{center}
{\renewcommand{\arraystretch}{2.4}\setlength{\tabcolsep}{5pt}%
\begin{tabular}{|c|c|c|c|c|c|c|c|c|c|}
\hline
angle & $0^{\circ}$ & $30^{\circ}$ & $45^{\circ}$ & $60^{\circ}$ & $90^{\circ}$ &
$120^{\circ}$ & $135^{\circ}$ & $150^{\circ}$ & $180^{\circ}$\\
\hline
en radians & $0$ & $\dfrac{\pi}{6}$ & $\dfrac{\pi}{4}$ & $\dfrac{\pi}{3}$ & $\dfrac{\pi}{2}$ &
$\dfrac{2\pi}{3}$ & $\dfrac{3\pi}{4}$ & $\dfrac{5\pi}{6}$ & $\pi$\\
\hline
cosinus & $1$ & $\dfrac{\sqrt{3}}{2}$ & $\dfrac{\sqrt{2}}{2}$ & $\dfrac{1}{2}$ & $0$ &
$-\dfrac{1}{2}$ & $-\dfrac{\sqrt{2}}{2}$ & $-\dfrac{\sqrt{3}}{2}$ & $-1$\\
\hline
\end{tabular}}
\end{center}
Deux angles qui se complètent à $180^{\circ}$ ont des cosinus opposés (points
symétriques par rapport à l'axe des ordonnées) : $\cos(180^{\circ}-a)=-\cos(a)$, par
exemple $\cos(120^{\circ})=-\cos(60^{\circ})$.
\end{remarquebox}

\begin{proprietebox}[Propriété 7 --- Signe, symétrie, orthogonalité]
\begin{itemize}
  \item $\vec{u}\cdot\vec{v}=\vec{v}\cdot\vec{u}$ (\textbf{symétrie}).
  \item Le produit scalaire de deux vecteurs non nuls est positif si leur angle est
        aigu, négatif s'il est obtus, nul s'il est droit.
  \item Deux vecteurs sont \textbf{orthogonaux} quand leur produit scalaire est nul :
        $\vec{u}\cdot\vec{v}=0$. Pour deux vecteurs non nuls, cela veut dire que leurs
        directions sont perpendiculaires ($\cos(90^{\circ})=0$) ; le vecteur nul est
        orthogonal à tout vecteur.
\end{itemize}
\end{proprietebox}

\begin{methodebox}[Méthode 9 --- Calculer un produit scalaire avec un angle]
Les deux vecteurs doivent avoir la même origine : sinon, on déplace l'un d'eux.
\etapeexemples
\begin{questions}[label=\alph*)]
  \item $\norme{\vec{u}}=3$, $\norme{\vec{v}}=4$ et l'angle de $\vec{u}$ et $\vec{v}$
        mesure $60^{\circ}$. Calculer $\vec{u}\cdot\vec{v}$.
  \item Dans le triangle $ABC$, $AB=5$, $AC=4$ et $\widehat{BAC}=135^{\circ}$. Calculer
        $\vect{AB}\cdot\vect{AC}$.
  \item $EFG$ est un triangle équilatéral de côté $6$. Calculer $\vect{EF}\cdot\vect{EG}$,
        puis $\vect{EF}\cdot\vect{FG}$.
\end{questions}
\begin{correction}
\cas{a)}
\begin{calculs}
\vec{u}\cdot\vec{v} &= \norme{\vec{u}}\times\norme{\vec{v}}\times\cos(60^{\circ}) &\qquad& \rem{on recopie la définition}\\
\vec{u}\cdot\vec{v} &= 3\times 4\times\dfrac{1}{2} && \rem{$\cos(60^{\circ})=\frac{1}{2}$}\\
\vec{u}\cdot\vec{v} &= \res{6} &&
\end{calculs}
\cas{b)}
\begin{calculs}
\vect{AB}\cdot\vect{AC} &= AB\times AC\times\cos\big(\widehat{BAC}\big) &\qquad& \rem{même origine $A$}\\
\vect{AB}\cdot\vect{AC} &= 5\times 4\times\left(-\dfrac{\sqrt{2}}{2}\right) && \rem{$\cos(135^{\circ})=-\frac{\sqrt{2}}{2}$}\\
\vect{AB}\cdot\vect{AC} &= \res{-10\sqrt{2}} && \rem{négatif : l'angle est obtus}
\end{calculs}
\cas{c)}\par\nopagebreak
\begin{minipage}[t]{0.58\linewidth}\raggedright
\begin{calculs}
\vect{EF}\cdot\vect{EG} &= EF\times EG\times\cos\big(\widehat{FEG}\big) &\qquad& \rem{même origine $E$}\\
\vect{EF}\cdot\vect{EG} &= 6\times 6\times\dfrac{1}{2} && \rem{angle de $60^{\circ}$}\\
\vect{EF}\cdot\vect{EG} &= \res{18} &&
\end{calculs}
\explique{$\vect{EF}$ et $\vect{FG}$ n'ont pas la même origine : on trace le représentant
$\vect{FK}$ de $\vect{EF}$ d'origine $F$.}{on prolonge $[EF]$ au-delà de $F$}
\begin{calculs}
\vect{EF}\cdot\vect{FG} &= \vect{FK}\cdot\vect{FG} &\qquad& \rem{même origine $F$}\\
\vect{EF}\cdot\vect{FG} &= 6\times 6\times\left(-\dfrac{1}{2}\right) && \rem{$\widehat{KFG}=180^{\circ}-60^{\circ}=120^{\circ}$}\\
\vect{EF}\cdot\vect{FG} &= \res{-18} &&
\end{calculs}
\end{minipage}\hfill
\begin{minipage}[t]{0.38\linewidth}
\vspace{0pt}\centering
\begin{tikzpicture}[scale=0.5]
  \draw[fig] (0,0) -- (3,0) -- (1.5,2.598) -- cycle;
  \draw[aide] (3,0) -- (6,0);
  \draw (3.8,0) arc (0:120:0.8);
  \node[font=\footnotesize] at (4.1,0.9) {$120^{\circ}$};
  \draw[vecteur] (3,0) -- node[midway,below,font=\footnotesize]{$\vect{FK}$} (6,0);
  \draw[vecteurB] (3,0) -- node[pos=0.75,right,font=\footnotesize]{$\vect{FG}$} (1.5,2.598);
  \pt{0,0}{below left}{E} \pt{3,0}{below}{F} \pt{1.5,2.598}{above}{G} \pt{6,0}{below}{K}
\end{tikzpicture}
\end{minipage}
\end{correction}

\etape{À vous de jouer}
\begin{questions}[label=\alph*)]
  \item $\norme{\vec{u}}=5$, $\norme{\vec{v}}=2$ et l'angle de $\vec{u}$ et $\vec{v}$
        mesure $60^{\circ}$. Calculer $\vec{u}\cdot\vec{v}$.
  \item Dans le triangle $RST$, $RS=3$, $RT=8$ et $\widehat{SRT}=150^{\circ}$. Calculer
        $\vect{RS}\cdot\vect{RT}$.
  \item $KLM$ est un triangle équilatéral de côté $4$. Calculer $\vect{KL}\cdot\vect{KM}$,
        puis $\vect{KL}\cdot\vect{LM}$.
\end{questions}
\reponse{a)~$5$ ;\; b)~$-12\sqrt{3}$ ;\; c)~$8$, puis $-8$.}
\end{methodebox}

\afaire{exercices 17 à 20}

\besoinplace{14\baselineskip}
\partiecours{Produit scalaire et projeté orthogonal}

\begin{definitionbox}[Définition 5 --- Projeté orthogonal d'un point]
\begin{minipage}[t]{0.58\linewidth}\raggedright
Le \textbf{projeté orthogonal} d'un point $M$ sur une droite $(d)$ est le point $H$ où la
perpendiculaire à $(d)$ qui passe par $M$ coupe $(d)$ : c'est l'« ombre » de $M$ sur
$(d)$ quand la lumière arrive perpendiculairement. Si $M$ est sur $(d)$, alors $H=M$.
\end{minipage}\hfill
\begin{minipage}[t]{0.38\linewidth}
\vspace{0pt}\centering
\begin{tikzpicture}[scale=0.7]
  \draw[fig] (-0.5,0) -- (4.5,0) node[right,font=\footnotesize]{$(d)$};
  \draw[aide] (2,2) -- (2,0);
  \draw[black!60] (2,0.25) -- (2.25,0.25) -- (2.25,0);
  \pt{2,2}{above}{M} \pt{2,0}{below}{H}
\end{tikzpicture}
\end{minipage}
\end{definitionbox}

\begin{proprietebox}[Théorème 3 --- Produit scalaire par projection]
Soit $A$ et $B$ deux points distincts, $C$ un point et $H$ son projeté orthogonal sur la
droite $(AB)$. Alors $\vect{AB}\cdot\vect{AC}=\vect{AB}\cdot\vect{AH}$, c'est-à-dire :
\begin{itemize}
  \item $\vect{AB}\cdot\vect{AC}=AB\times AH$ si $\vect{AB}$ et $\vect{AH}$ ont le
        \textbf{même sens} ;
  \item $\vect{AB}\cdot\vect{AC}=-AB\times AH$ s'ils sont de \textbf{sens contraires}.
\end{itemize}
\begin{center}
\begin{tikzpicture}[scale=0.65]
  \draw[aide] (2.5,2) -- (2.5,0);
  \draw[black!60] (2.5,0.25) -- (2.25,0.25) -- (2.25,0);
  \draw[vecteur] (0,0) -- node[pos=0.85,above,font=\footnotesize]{$\vect{AB}$} (4,0);
  \draw[vecteurB] (0,0) -- node[midway,above left,font=\footnotesize]{$\vect{AC}$} (2.5,2);
  \pt{0,0}{left}{A} \pt{4,0}{right}{B} \pt{2.5,2}{above}{C} \pt{2.5,0}{below}{H}
  \node[font=\footnotesize] at (2,-1.3) {même sens : $AB\times AH$};
\end{tikzpicture}
\qquad\qquad
\begin{tikzpicture}[scale=0.65]
  \draw[aide] (-2,0) -- (0,0);
  \draw[aide] (-2,2) -- (-2,0);
  \draw[black!60] (-2,0.25) -- (-1.75,0.25) -- (-1.75,0);
  \draw[vecteur] (0,0) -- node[midway,below,font=\footnotesize]{$\vect{AB}$} (3,0);
  \draw[vecteurB] (0,0) -- node[midway,above right,font=\footnotesize]{$\vect{AC}$} (-2,2);
  \pt{0,0}{below}{A} \pt{3,0}{right}{B} \pt{-2,2}{above}{C} \pt{-2,0}{below}{H}
  \node[font=\footnotesize] at (0.5,-1.3) {sens contraires : $-AB\times AH$};
\end{tikzpicture}
\end{center}
En particulier, si $C$ est sur la droite $(AB)$ : $\vect{AB}\cdot\vect{AC}=AB\times AC$
quand $\vect{AB}$ et $\vect{AC}$ ont le même sens, $-AB\times AC$ sinon.

\smallskip
\demo{pour un angle $\widehat{BAC}$ aigu, dans le triangle $AHC$ rectangle en $H$ :}
\begin{calculs}
\cos\big(\widehat{BAC}\big) &= \dfrac{AH}{AC} &\qquad& \rem{côté adjacent sur hypoténuse}\\
\vect{AB}\cdot\vect{AC} &= AB\times AC\times\dfrac{AH}{AC} && \rem{définition 4}\\
\vect{AB}\cdot\vect{AC} &= AB\times AH && \rem{on simplifie par $AC$}
\end{calculs}
Pour un angle $\widehat{BAC}$ obtus, $H$ est de l'autre côté de $A$ et
$\widehat{HAC}=180^{\circ}-\widehat{BAC}$ est aigu :
\begin{calculs}
\cos\big(\widehat{BAC}\big) &= -\cos\big(\widehat{HAC}\big) &\qquad& \rem{angles qui se complètent à $180^{\circ}$}\\
\cos\big(\widehat{BAC}\big) &= -\dfrac{AH}{AC} && \rem{dans le triangle $AHC$ rectangle en $H$}\\
\vect{AB}\cdot\vect{AC} &= -AB\times AH && \rem{définition 4, on simplifie par $AC$}
\end{calculs}
\end{proprietebox}

\begin{methodebox}[Méthode 10 --- Calculer un produit scalaire par projection]
On projette l'extrémité du second vecteur sur la droite qui porte le premier, puis on
compare les sens.
\etapeexemples
\begin{minipage}[t]{0.62\linewidth}\raggedright
$ABCD$ est un carré de côté $4$ et $I$ est le milieu de $[CD]$. Calculer :
\begin{questions}[label=\alph*)]
  \item $\vect{AB}\cdot\vect{AC}$ ;
  \item $\vect{AB}\cdot\vect{AD}$ ;
  \item $\vect{AB}\cdot\vect{CD}$ ;
  \item $\vect{AB}\cdot\vect{AI}$.
\end{questions}
\end{minipage}\hfill
\begin{minipage}[t]{0.34\linewidth}
\vspace{0pt}\centering
\carremilieu{A}{B}{C}{D}{I}
\end{minipage}

\smallskip
\begin{minipage}[t]{0.62\linewidth}\raggedright
\begin{questions}[label=\alph*),start=5]
  \item Dans le triangle $EFG$ ci-contre, $H$ est le projeté orthogonal de $G$ sur la
        droite $(EF)$, avec $EF=5$ et $EH=3$. Calculer $\vect{EF}\cdot\vect{EG}$.
\end{questions}
\end{minipage}\hfill
\begin{minipage}[t]{0.34\linewidth}
\vspace{0pt}\centering
\begin{tikzpicture}[scale=0.4]
  \draw[aide] (-3,0) -- (0,0);
  \draw[aide] (-3,2.5) -- (-3,0);
  \draw[black!60] (-3,0.4) -- (-2.6,0.4) -- (-2.6,0);
  \draw[fig] (0,0) -- (5,0) -- (-3,2.5) -- cycle;
  \pt{0,0}{below}{E} \pt{5,0}{below right}{F} \pt{-3,2.5}{above}{G} \pt{-3,0}{below}{H}
\end{tikzpicture}
\end{minipage}
\begin{correction}
\cas{a)}
\begin{calculs}
\vect{AB}\cdot\vect{AC} &= \vect{AB}\cdot\vect{AB} &\qquad& \rem{$B$ est le projeté de $C$ sur $(AB)$}\\
\vect{AB}\cdot\vect{AC} &= AB^{2} && \rem{carré scalaire}\\
\vect{AB}\cdot\vect{AC} &= \res{16} && \rem{$4^{2}$}
\end{calculs}
\cas{b)}
\explique{$\vect{AB}\cdot\vect{AD}=\res{0}$}{$(AB)$ et $(AD)$ sont perpendiculaires :
le projeté de $D$ est $A$}
\cas{c)}
\begin{calculs}
\vect{AB}\cdot\vect{CD} &= \vect{AB}\cdot\vect{BA} &\qquad& \rem{$\vect{CD}=\vect{BA}$ : on se ramène à l'origine $A$ ou $B$}\\
\vect{AB}\cdot\vect{CD} &= -AB\times AB && \rem{$\vect{AB}$ et $\vect{BA}$ colinéaires, de sens contraires}\\
\vect{AB}\cdot\vect{CD} &= \res{-16} &&
\end{calculs}
\cas{d)}
\explique{Le projeté de $I$ sur $(AB)$ est le milieu $J$ de $[AB]$, et $AJ=2$.}{$I$ est au
milieu du côté du haut}
\begin{calculs}
\vect{AB}\cdot\vect{AI} &= AB\times AJ &\qquad& \rem{$\vect{AB}$ et $\vect{AJ}$ ont le même sens}\\
\vect{AB}\cdot\vect{AI} &= 4\times 2 &&\\
\vect{AB}\cdot\vect{AI} &= \res{8} &&
\end{calculs}
\cas{e)}
\begin{calculs}
\vect{EF}\cdot\vect{EG} &= -EF\times EH &\qquad& \rem{$\vect{EH}$ va dans le sens contraire de $\vect{EF}$}\\
\vect{EF}\cdot\vect{EG} &= -5\times 3 &&\\
\vect{EF}\cdot\vect{EG} &= \res{-15} &&
\end{calculs}
\end{correction}

\etape{À vous de jouer}
\begin{minipage}[t]{0.62\linewidth}\raggedright
$KLMN$ est un carré de côté $6$ et $J$ est le milieu de $[MN]$. Calculer :
\begin{questions}[label=\alph*)]
  \item $\vect{KL}\cdot\vect{KM}$ ;
  \item $\vect{KL}\cdot\vect{KN}$ ;
  \item $\vect{KL}\cdot\vect{MN}$ ;
  \item $\vect{KL}\cdot\vect{KJ}$.
\end{questions}
\end{minipage}\hfill
\begin{minipage}[t]{0.34\linewidth}
\vspace{0pt}\centering
\carremilieu{K}{L}{M}{N}{J}
\end{minipage}

\smallskip
\begin{minipage}[t]{0.62\linewidth}\raggedright
\begin{questions}[label=\alph*),start=5]
  \item Dans le triangle $RST$ ci-contre, $H$ est le projeté orthogonal de $T$ sur la
        droite $(RS)$, avec $RS=4$ et $RH=2$. Calculer $\vect{RS}\cdot\vect{RT}$.
\end{questions}
\end{minipage}\hfill
\begin{minipage}[t]{0.34\linewidth}
\vspace{0pt}\centering
\begin{tikzpicture}[scale=0.45]
  \draw[aide] (-2,0) -- (0,0);
  \draw[aide] (-2,2.5) -- (-2,0);
  \draw[black!60] (-2,0.4) -- (-1.6,0.4) -- (-1.6,0);
  \draw[fig] (0,0) -- (4,0) -- (-2,2.5) -- cycle;
  \pt{0,0}{below}{R} \pt{4,0}{below right}{S} \pt{-2,2.5}{above}{T} \pt{-2,0}{below}{H}
\end{tikzpicture}
\end{minipage}
\reponse{a)~$36$ ;\; b)~$0$ ;\; c)~$-36$ ;\; d)~$18$ ;\; e)~$-8$.}
\end{methodebox}

\afaire{exercices 21 à 23}

\partiecours{Règles de calcul}

Avec le produit scalaire, on développe comme en calcul littéral, où
\begin{center}
$(a+b)(c+d)=ac+ad+bc+bd$ \qquad $(a+b)^{2}=a^{2}+2ab+b^{2}$
\end{center}

\begin{proprietebox}[Propriété 8 --- Règles de calcul (admises)]
Pour tous vecteurs $\vec{u}$, $\vec{v}$, $\vec{w}$ et tout nombre réel $k$ :
\begin{center}
$\vec{u}\cdot(\vec{v}+\vec{w})=\vec{u}\cdot\vec{v}+\vec{u}\cdot\vec{w}$
\qquad
$(\vec{u}+\vec{v})\cdot\vec{w}=\vec{u}\cdot\vec{w}+\vec{v}\cdot\vec{w}$
\qquad
$(k\vec{u})\cdot\vec{v}=\vec{u}\cdot(k\vec{v})=k\,(\vec{u}\cdot\vec{v})$
\end{center}
On dit que le produit scalaire est \textbf{bilinéaire} : on développe, et on sort les
nombres réels.
\end{proprietebox}

\begin{proprietebox}[Propriété 9 --- Identités remarquables]
Pour tous vecteurs $\vec{u}$ et $\vec{v}$ :
\begin{center}
$(\vec{u}+\vec{v})^{2}=\vec{u}^{2}+2\,\vec{u}\cdot\vec{v}+\vec{v}^{2}$
\qquad
$(\vec{u}-\vec{v})^{2}=\vec{u}^{2}-2\,\vec{u}\cdot\vec{v}+\vec{v}^{2}$
\qquad
$(\vec{u}+\vec{v})\cdot(\vec{u}-\vec{v})=\vec{u}^{2}-\vec{v}^{2}$
\end{center}
Avec les normes : $\norme{\vec{u}+\vec{v}}^{2}=\norme{\vec{u}}^{2}+2\,\vec{u}\cdot\vec{v}+\norme{\vec{v}}^{2}$.

\smallskip
\demo{de la première (les deux autres se démontrent de la même façon) :}
\begin{calculs}
(\vec{u}+\vec{v})^{2} &= (\vec{u}+\vec{v})\cdot(\vec{u}+\vec{v}) &\qquad& \rem{carré scalaire}\\
(\vec{u}+\vec{v})^{2} &= \vec{u}\cdot\vec{u}+\vec{u}\cdot\vec{v}+\vec{v}\cdot\vec{u}+\vec{v}\cdot\vec{v} && \rem{on développe (propriété 8)}\\
(\vec{u}+\vec{v})^{2} &= \vec{u}^{2}+2\,\vec{u}\cdot\vec{v}+\vec{v}^{2} && \rem{symétrie : $\vec{v}\cdot\vec{u}=\vec{u}\cdot\vec{v}$}
\end{calculs}
\end{proprietebox}

\begin{methodebox}[Méthode 11 --- Calculer avec les règles du produit scalaire]
\etapeexemples
On sait que $\norme{\vec{u}}=3$, $\norme{\vec{v}}=2$ et $\vec{u}\cdot\vec{v}=-1$. Calculer :
\begin{questions}[label=\alph*)]
  \item $(2\vec{u})\cdot(3\vec{v})$ ;
  \item $\vec{u}\cdot(\vec{u}+\vec{v})$ ;
  \item $\norme{\vec{u}+\vec{v}}^{2}$, puis $\norme{\vec{u}+\vec{v}}$ ;
  \item $(\vec{u}-2\vec{v})\cdot(\vec{u}+\vec{v})$.
\end{questions}
\begin{correction}
\cas{a)}
\begin{calculs}
(2\vec{u})\cdot(3\vec{v}) &= 2\times 3\times(\vec{u}\cdot\vec{v}) &\qquad& \rem{on sort les nombres réels}\\
(2\vec{u})\cdot(3\vec{v}) &= 6\times(-1) &&\\
(2\vec{u})\cdot(3\vec{v}) &= \res{-6} &&
\end{calculs}
\cas{b)}
\begin{calculs}
\vec{u}\cdot(\vec{u}+\vec{v}) &= \vec{u}\cdot\vec{u}+\vec{u}\cdot\vec{v} &\qquad& \rem{on développe}\\
\vec{u}\cdot(\vec{u}+\vec{v}) &= \norme{\vec{u}}^{2}+\vec{u}\cdot\vec{v} && \rem{carré scalaire}\\
\vec{u}\cdot(\vec{u}+\vec{v}) &= 9+(-1) && \rem{$3^{2}=9$}\\
\vec{u}\cdot(\vec{u}+\vec{v}) &= \res{8} &&
\end{calculs}
\cas{c)}
\begin{calculs}
\norme{\vec{u}+\vec{v}}^{2} &= \norme{\vec{u}}^{2}+2\,\vec{u}\cdot\vec{v}+\norme{\vec{v}}^{2} &\qquad& \rem{identité remarquable}\\
\norme{\vec{u}+\vec{v}}^{2} &= 9+2\times(-1)+4 &&\\
\norme{\vec{u}+\vec{v}}^{2} &= \res{11} &&\\
\norme{\vec{u}+\vec{v}} &= \res{\sqrt{11}} && \rem{une norme est positive}
\end{calculs}
\cas{d)}
\begin{calculs}
(\vec{u}-2\vec{v})\cdot(\vec{u}+\vec{v}) &= \vec{u}\cdot\vec{u}+\vec{u}\cdot\vec{v}-2\,\vec{v}\cdot\vec{u}-2\,\vec{v}\cdot\vec{v} &\qquad& \rem{on développe}\\
(\vec{u}-2\vec{v})\cdot(\vec{u}+\vec{v}) &= \norme{\vec{u}}^{2}-\vec{u}\cdot\vec{v}-2\norme{\vec{v}}^{2} && \rem{$\vec{u}\cdot\vec{v}-2\,\vec{u}\cdot\vec{v}=-\vec{u}\cdot\vec{v}$}\\
(\vec{u}-2\vec{v})\cdot(\vec{u}+\vec{v}) &= 9-(-1)-2\times 4 &&\\
(\vec{u}-2\vec{v})\cdot(\vec{u}+\vec{v}) &= \res{2} &&
\end{calculs}
\end{correction}

\etape{À vous de jouer}
On sait que $\norme{\vec{u}}=4$, $\norme{\vec{v}}=1$ et $\vec{u}\cdot\vec{v}=2$. Calculer :
\begin{questions}[label=\alph*)]
  \item $(3\vec{u})\cdot(2\vec{v})$ ;
  \item $\vec{u}\cdot(\vec{u}+\vec{v})$ ;
  \item $\norme{\vec{u}+\vec{v}}^{2}$, puis $\norme{\vec{u}+\vec{v}}$ ;
  \item $(\vec{u}-2\vec{v})\cdot(\vec{u}+\vec{v})$.
\end{questions}
\reponse{a)~$12$ ;\; b)~$18$ ;\; c)~$21$, puis $\sqrt{21}$ ;\; d)~$12$.}
\end{methodebox}

\afaire{exercices 24 à 26}

\partiecours{Produit scalaire dans un repère orthonormé}

\begin{proprietebox}[Théorème 4 --- Expression avec les coordonnées]
Dans un repère \textbf{orthonormé}, pour $\vec{u}\coord{x}{y}$ et $\vec{v}\coord{x'}{y'}$ :
\[
  \vec{u}\cdot\vec{v}=xx'+yy'
  \qquad\text{(abscisses multipliées, plus ordonnées multipliées).}
\]
On retrouve $\vec{u}\cdot\vec{u}=x^{2}+y^{2}=\norme{\vec{u}}^{2}$.

\smallskip
\demo{avec $\vec{u}=x\vec{\imath}+y\vec{\jmath}$ et $\vec{v}=x'\vec{\imath}+y'\vec{\jmath}$ :}
\begin{calculs}
\vec{u}\cdot\vec{v} &= (x\vec{\imath}+y\vec{\jmath})\cdot(x'\vec{\imath}+y'\vec{\jmath}) &\qquad& \rem{définition 1}\\
\vec{u}\cdot\vec{v} &= xx'\,\vec{\imath}\cdot\vec{\imath}+xy'\,\vec{\imath}\cdot\vec{\jmath}+yx'\,\vec{\jmath}\cdot\vec{\imath}+yy'\,\vec{\jmath}\cdot\vec{\jmath} && \rem{on développe (propriété 8)}\\
\vec{u}\cdot\vec{v} &= xx'\times 1+xy'\times 0+yx'\times 0+yy'\times 1 && \rem{$\vec{\imath}^{2}=\norme{\vec{\imath}}^{2}=1$ ; $\vec{\imath}$ et $\vec{\jmath}$ orthogonaux}\\
\vec{u}\cdot\vec{v} &= xx'+yy' &&
\end{calculs}

\smallskip
\begin{itemize}
  \item $\vec{u}$ et $\vec{v}$ sont \textbf{orthogonaux} quand $xx'+yy'=0$.
  \item Avec $\vec{\imath}\coord{1}{0}$ et $\vec{\jmath}\coord{0}{1}$ :
        $\vec{u}\cdot\vec{\imath}=x$ et $\vec{u}\cdot\vec{\jmath}=y$. Les coordonnées
        de $\vec{u}$ sont ses produits scalaires avec les vecteurs de la base.
\end{itemize}
À ne pas confondre : $xy'-yx'=0$ (déterminant) pour la \textbf{colinéarité},
$xx'+yy'=0$ (produit scalaire) pour l'\textbf{orthogonalité}.
\end{proprietebox}

\begin{methodebox}[Méthode 12 --- Calculer un produit scalaire avec les coordonnées, vecteurs orthogonaux]
\etapeexemples
Le repère est orthonormé.
\begin{questions}[label=\alph*)]
  \item Calculer $\vec{u}\cdot\vec{v}$ pour $\vec{u}\coord{4}{-3}$ et $\vec{v}\coord{2}{5}$.
  \item Les vecteurs $\vec{a}\coord{6}{-4}$ et $\vec{b}\coord{2}{3}$ sont-ils orthogonaux ?
  \item On considère les points $A(1\,;-1)$, $B(5\,;1)$ et $C(0\,;1)$. Le triangle $ABC$
        est-il rectangle en $A$ ?
  \item Calculer le nombre réel $m$ pour lequel les vecteurs $\vec{r}\coord{m}{3}$ et
        $\vec{s}\coord{6}{-4}$ sont orthogonaux.
\end{questions}
\begin{correction}
\cas{a)}
\begin{calculs}
\vec{u}\cdot\vec{v} &= xx'+yy' &\qquad& \rem{on recopie la formule}\\
\vec{u}\cdot\vec{v} &= 4\times 2+(-3)\times 5 && \rem{abscisses ensemble, ordonnées ensemble}\\
\vec{u}\cdot\vec{v} &= 8-15 &&\\
\vec{u}\cdot\vec{v} &= \res{-7} &&
\end{calculs}
\cas{b)}
\begin{calculs}
\vec{a}\cdot\vec{b} &= 6\times 2+(-4)\times 3 &\qquad& \rem{$xx'+yy'$}\\
\vec{a}\cdot\vec{b} &= 0 && \rem{$12-12$}
\end{calculs}
\explique{Le produit scalaire est nul : $\vec{a}$ et $\vec{b}$ \res{sont orthogonaux}.}{théorème 4}
\cas{c)}
\begin{calculs}
\vect{AB} & \coord{4}{2} &\qquad& \rem{les deux côtés issus de $A$ : $5-1$ et $1-(-1)$}\\
\vect{AC} & \coord{-1}{2} && \rem{$0-1$ et $1-(-1)$}\\
\vect{AB}\cdot\vect{AC} &= 4\times(-1)+2\times 2 && \rem{$xx'+yy'$}\\
\vect{AB}\cdot\vect{AC} &= 0 &&
\end{calculs}
\explique{$\vect{AB}$ et $\vect{AC}$ sont orthogonaux : le triangle $ABC$ \res{est rectangle
en $A$}.}{les droites $(AB)$ et $(AC)$ sont perpendiculaires}
\cas{d)}
\begin{calculs}
\vec{r}\cdot\vec{s} &= m\times 6+3\times(-4) &\qquad& \rem{$xx'+yy'$}\\
\vec{r}\cdot\vec{s} &= 6m-12 &&\\
6m-12 &= 0 && \rem{le produit scalaire doit être nul}\\
m &= \res{2} && \rem{on ajoute $12$, puis on divise par $6$}
\end{calculs}
\end{correction}

\etape{À vous de jouer}
Le repère est orthonormé.
\begin{questions}[label=\alph*)]
  \item Calculer $\vec{u}\cdot\vec{v}$ pour $\vec{u}\coord{-2}{5}$ et $\vec{v}\coord{3}{1}$.
  \item Les vecteurs $\vec{a}\coord{5}{2}$ et $\vec{b}\coord{-4}{10}$ sont-ils orthogonaux ?
  \item On considère les points $E(-1\,;2)$, $F(3\,;4)$ et $G(1\,;-2)$. Le triangle $EFG$
        est-il rectangle en $E$ ?
  \item Calculer le nombre réel $m$ pour lequel les vecteurs $\vec{r}\coord{m}{-2}$ et
        $\vec{s}\coord{5}{10}$ sont orthogonaux.
\end{questions}
\reponse{a)~$-1$ ;\; b)~oui ;\; c)~oui ;\; d)~$m=4$.}
\end{methodebox}

\afaire{exercices 27 à 30}

\partiecours{Produit scalaire et normes}

Les identités remarquables (propriété 9) relient le produit scalaire aux longueurs :
quand on connaît les trois côtés d'un triangle, on calcule un produit scalaire sans angle
ni projeté.

\begin{proprietebox}[Propriété 10 --- Produit scalaire avec des normes]
\begin{minipage}[t]{0.58\linewidth}\raggedright
Pour tous vecteurs $\vec{u}$ et $\vec{v}$ :
\[
  \vec{u}\cdot\vec{v}=\frac{1}{2}\left(\norme{\vec{u}+\vec{v}}^{2}-\norme{\vec{u}}^{2}-\norme{\vec{v}}^{2}\right).
\]
Dans un triangle $ABC$ :
\[
  \vect{AB}\cdot\vect{AC}=\frac{1}{2}\left(AB^{2}+AC^{2}-BC^{2}\right).
\]
Le côté retranché, $BC$, est celui qui \textbf{fait face} au sommet $A$ commun aux deux
vecteurs.
\end{minipage}\hfill
\begin{minipage}[t]{0.38\linewidth}
\vspace{0pt}\centering
\begin{tikzpicture}[scale=0.6]
  \draw[aide] (5,0) -- (1.5,2.6);
  \node[font=\footnotesize,above right] at (3.25,1.3) {$BC$};
  \draw[vecteur] (0,0) -- node[midway,below,font=\footnotesize]{$\vect{AB}$} (5,0);
  \draw[vecteurB] (0,0) -- node[midway,above left,font=\footnotesize]{$\vect{AC}$} (1.5,2.6);
  \pt{0,0}{below left}{A} \pt{5,0}{below right}{B} \pt{1.5,2.6}{above}{C}
\end{tikzpicture}
\end{minipage}

\smallskip
\demo{dans le triangle, avec la relation de Chasles et la propriété 9 :}
\begin{calculs}
BC^{2} &= \norme{\vect{AC}-\vect{AB}}^{2} &\qquad& \rem{$\vect{BC}=\vect{BA}+\vect{AC}=\vect{AC}-\vect{AB}$}\\
BC^{2} &= AC^{2}-2\,\vect{AB}\cdot\vect{AC}+AB^{2} && \rem{identité remarquable}\\
2\,\vect{AB}\cdot\vect{AC} &= AB^{2}+AC^{2}-BC^{2} && \rem{on isole le produit scalaire}\\
\vect{AB}\cdot\vect{AC} &= \frac{1}{2}\left(AB^{2}+AC^{2}-BC^{2}\right) && \rem{on divise par $2$}
\end{calculs}
La formule avec $\norme{\vec{u}+\vec{v}}$ s'obtient de même, en isolant $\vec{u}\cdot\vec{v}$
dans $\norme{\vec{u}+\vec{v}}^{2}=\norme{\vec{u}}^{2}+2\,\vec{u}\cdot\vec{v}+\norme{\vec{v}}^{2}$.
\end{proprietebox}

\begin{methodebox}[Méthode 13 --- Calculer un produit scalaire avec des normes]
On repère le sommet commun aux deux vecteurs ; on retranche le carré du côté qui lui
fait face.
\etapeexemples
\begin{questions}[label=\alph*)]
  \item $\norme{\vec{u}}=3$, $\norme{\vec{v}}=5$ et $\norme{\vec{u}+\vec{v}}=7$. Calculer
        $\vec{u}\cdot\vec{v}$.
  \item Dans le triangle $ABC$, $AB=4$, $AC=6$ et $BC=5$. Calculer $\vect{AB}\cdot\vect{AC}$.
  \item Dans le même triangle $ABC$, calculer $\vect{BA}\cdot\vect{BC}$.
\end{questions}
\begin{correction}
\cas{a)}
\begin{calculs}
\vec{u}\cdot\vec{v} &= \frac{1}{2}\left(\norme{\vec{u}+\vec{v}}^{2}-\norme{\vec{u}}^{2}-\norme{\vec{v}}^{2}\right) &\qquad& \rem{on recopie la formule}\\
\vec{u}\cdot\vec{v} &= \frac{1}{2}\left(49-9-25\right) && \rem{$7^{2}$, $3^{2}$ et $5^{2}$}\\
\vec{u}\cdot\vec{v} &= \res{\frac{15}{2}} &&
\end{calculs}
\cas{b)}
\begin{calculs}
\vect{AB}\cdot\vect{AC} &= \frac{1}{2}\left(AB^{2}+AC^{2}-BC^{2}\right) &\qquad& \rem{sommet commun $A$ : on retranche $BC^{2}$}\\
\vect{AB}\cdot\vect{AC} &= \frac{1}{2}\left(16+36-25\right) &&\\
\vect{AB}\cdot\vect{AC} &= \res{\frac{27}{2}} &&
\end{calculs}
\cas{c)}
\begin{calculs}
\vect{BA}\cdot\vect{BC} &= \frac{1}{2}\left(BA^{2}+BC^{2}-AC^{2}\right) &\qquad& \rem{sommet commun $B$ : on retranche $AC^{2}$}\\
\vect{BA}\cdot\vect{BC} &= \frac{1}{2}\left(16+25-36\right) &&\\
\vect{BA}\cdot\vect{BC} &= \res{\frac{5}{2}} &&
\end{calculs}
\end{correction}

\etape{À vous de jouer}
\begin{questions}[label=\alph*)]
  \item $\norme{\vec{u}}=2$, $\norme{\vec{v}}=4$ et $\norme{\vec{u}+\vec{v}}=5$. Calculer
        $\vec{u}\cdot\vec{v}$.
  \item Dans le triangle $EFG$, $EF=5$, $EG=7$ et $FG=6$. Calculer $\vect{EF}\cdot\vect{EG}$.
  \item Dans le même triangle $EFG$, calculer $\vect{FE}\cdot\vect{FG}$.
\end{questions}
\reponse{a)~$\frac{5}{2}$ ;\; b)~$19$ ;\; c)~$6$.}
\end{methodebox}

\afaire{exercices 31 à 33}

\partiecours{Vecteur normal et équation cartésienne d'une droite}

Rappel de seconde (chapitre « Équations cartésiennes de droites ») : une droite a une
équation de la forme $ax+by+c=0$, et $\vec{u}\coord{-b}{a}$ en est un vecteur directeur.
Le produit scalaire donne un autre vecteur, perpendiculaire à la droite, qu'on lit tout
aussi vite.

\begin{definitionbox}[Définition 6 --- Vecteur normal]
\begin{minipage}[t]{0.58\linewidth}\raggedright
Un \textbf{vecteur normal} à une droite $(d)$ est un vecteur $\vec{n}$ non nul orthogonal à
un vecteur directeur $\vec{u}$ de $(d)$ : sa direction est perpendiculaire à $(d)$.
\end{minipage}\hfill
\begin{minipage}[t]{0.38\linewidth}
\vspace{0pt}\centering
\begin{tikzpicture}[scale=0.8]
  \draw[fig] (-1,-0.5) -- (4,2) node[right,font=\footnotesize]{$(d)$};
  \draw[vecteur] (0,0) -- node[midway,below right,font=\footnotesize]{$\vec{u}$} (2.4,1.2);
  \draw[vecteurB] (0,0) -- node[midway,left,font=\footnotesize]{$\vec{n}$} (-0.8,1.6);
  \draw[black!60] (0.27,0.13) -- (0.14,0.4) -- (-0.13,0.27);
  \pt{0,0}{below}{A}
\end{tikzpicture}
\end{minipage}
\end{definitionbox}

\begin{proprietebox}[Théorème 5 --- Vecteur normal et équation cartésienne]
Le repère est orthonormé.
\begin{itemize}
  \item La droite d'équation $ax+by+c=0$ a pour vecteur normal $\vec{n}\coord{a}{b}$.
  \item Réciproquement, la droite qui passe par un point $A$ et de vecteur normal
        $\vec{n}\coord{a}{b}$ a une équation de la forme $ax+by+c=0$ : un point $M$ est sur
        cette droite quand $\vect{AM}\cdot\vec{n}=0$.
\end{itemize}
\demo{}
\begin{calculs}
\vec{u}\cdot\vec{n} &= (-b)\times a+a\times b &\qquad& \rem{$\vec{u}\coord{-b}{a}$ directeur, $xx'+yy'$}\\
\vec{u}\cdot\vec{n} &= 0 && \rem{$\vec{n}$ est orthogonal à $\vec{u}$}
\end{calculs}
Pour la réciproque, avec $A(x_{A}\,;y_{A})$ et $M(x\,;y)$ : $M$ est sur la droite quand
$\vect{AM}$ est orthogonal à $\vec{n}$.
\begin{calculs}
\vect{AM}\cdot\vec{n} &= a(x-x_{A})+b(y-y_{A}) &\qquad& \rem{$\vect{AM}\coord{x-x_{A}}{y-y_{A}}$}\\
\vect{AM}\cdot\vec{n} &= ax+by+c && \rem{en posant $c=-ax_{A}-by_{A}$}
\end{calculs}
\end{proprietebox}

\begin{methodebox}[Méthode 14 --- Trouver une équation cartésienne avec un vecteur normal]
On lit $a$ et $b$ sur le vecteur normal ; on trouve $c$ en écrivant que le point connu
vérifie l'équation.
\etapeexemples
Le repère est orthonormé.
\begin{questions}[label=\alph*)]
  \item Lire un vecteur normal et un vecteur directeur de la droite d'équation
        $3x-2y+5=0$.
  \item Trouver une équation cartésienne de la droite $(d)$ qui passe par $A(2\,;-1)$, de
        vecteur normal $\vec{n}\coord{3}{4}$.
  \item On considère les points $B(1\,;2)$, $C(-1\,;-2)$ et $D(4\,;0)$. Trouver une équation
        cartésienne de la hauteur issue de $B$ dans le triangle $BCD$.
  \item On considère les points $E(-1\,;3)$ et $F(3\,;1)$. Trouver une équation cartésienne
        de la médiatrice de $[EF]$.
\end{questions}
\begin{correction}
\cas{a)}
\explique{$3x+(-2)y+5=0$ : $a=3$ et $b=-2$, donc $\res{\vec{n}\coord{3}{-2}}$ est normal.}{on
lit $a$ et $b$ avec leurs signes}
\explique{Le vecteur $\vec{u}\coord{-b}{a}$, soit $\res{\vec{u}\coord{2}{3}}$, est
directeur.}{on vérifie : $\vec{n}\cdot\vec{u}=6-6=0$}
\cas{b)}
\begin{calculs}
3x+4y+c &= 0 &\qquad& \rem{$a=3$ et $b=4$, lus sur $\vec{n}$}\\
3\times 2+4\times(-1)+c &= 0 && \rem{$A(2\,;-1)$ est sur $(d)$}\\
2+c &= 0 &&\\
c &= -2 &&
\end{calculs}
\explique{Une équation de $(d)$ est $\res{3x+4y-2=0}$.}{on remplace $c$}
\cas{c)}
\explique{La hauteur issue de $B$ passe par $B$ et elle est perpendiculaire à $(CD)$ :
$\vect{CD}$ est un vecteur normal.}{on cherche d'abord le vecteur normal}
\begin{calculs}
\vect{CD} & \coord{5}{2} &\qquad& \rem{$4-(-1)$ et $0-(-2)$}\\
5x+2y+c &= 0 && \rem{$a=5$ et $b=2$}\\
5\times 1+2\times 2+c &= 0 && \rem{$B(1\,;2)$ est sur la hauteur}\\
c &= -9 && \rem{$9+c=0$}
\end{calculs}
\explique{Une équation de la hauteur est $\res{5x+2y-9=0}$.}{on remplace $c$}
\cas{d)}
\explique{La médiatrice de $[EF]$ passe par le milieu $I$ de $[EF]$ et elle est
perpendiculaire à $(EF)$ : $\vect{EF}$ est un vecteur normal.}{définition de la
médiatrice}
\begin{calculs}
I & \left(\dfrac{-1+3}{2}\,;\dfrac{3+1}{2}\right) &\qquad& \rem{coordonnées du milieu}\\
I & (1\,;2) &&\\
\vect{EF} & \coord{4}{-2} && \rem{$3-(-1)$ et $1-3$}\\
4x-2y+c &= 0 && \rem{$a=4$ et $b=-2$}\\
4\times 1-2\times 2+c &= 0 && \rem{$I(1\,;2)$ est sur la médiatrice}\\
c &= 0 &&
\end{calculs}
\explique{Une équation de la médiatrice est $\res{4x-2y=0}$, ou encore $2x-y=0$.}{on peut
diviser les deux membres par $2$}
\end{correction}

\etape{À vous de jouer}
Le repère est orthonormé.
\begin{questions}[label=\alph*)]
  \item Lire un vecteur normal et un vecteur directeur de la droite d'équation
        $-x+4y-3=0$.
  \item Trouver une équation cartésienne de la droite $(d)$ qui passe par $A(-3\,;1)$, de
        vecteur normal $\vec{n}\coord{2}{-5}$.
  \item On considère les points $K(2\,;3)$, $L(0\,;-1)$ et $M(3\,;1)$. Trouver une équation
        cartésienne de la hauteur issue de $K$ dans le triangle $KLM$.
  \item On considère les points $R(-2\,;-1)$ et $S(4\,;3)$. Trouver une équation cartésienne
        de la médiatrice de $[RS]$.
\end{questions}
\reponse{a)~$\vec{n}\coord{-1}{4}$ et $\vec{u}\coord{-4}{-1}$ ;\; b)~$2x-5y+11=0$ ;\;
c)~$3x+2y-12=0$ ;\; d)~$6x+4y-10=0$, ou $3x+2y-5=0$.}
\end{methodebox}

\begin{proprietebox}[Propriété 11 --- Droites parallèles, droites perpendiculaires]
Le repère est orthonormé. Deux droites $(d)$ et $(d')$ ont pour vecteurs normaux
$\vec{n}$ et $\vec{n}\,{}'$.
\begin{itemize}
  \item $(d)$ et $(d')$ sont \textbf{perpendiculaires} quand $\vec{n}\cdot\vec{n}\,{}'=0$ ;
  \item $(d)$ et $(d')$ sont \textbf{parallèles} quand $\vec{n}$ et $\vec{n}\,{}'$ sont
        colinéaires : $\det(\vec{n}\,;\vec{n}\,{}')=0$ (théorème 1).
\end{itemize}
En effet, deux droites sont perpendiculaires quand leurs perpendiculaires le sont, et
parallèles quand leurs perpendiculaires le sont.
\end{proprietebox}

\begin{methodebox}[Méthode 15 --- Démontrer que deux droites sont parallèles ou perpendiculaires]
On lit un vecteur normal sur chaque équation (cartésienne), puis on calcule leur produit
scalaire et leur déterminant.
\etapeexemples
Le repère est orthonormé. Pour chaque couple de droites, dire si elles sont parallèles,
perpendiculaires, ou ni l'un ni l'autre.
\begin{questions}[label=\alph*)]
  \item $(d_{1})$ : $2x-3y+1=0$ et $(d_{2})$ : $3x+2y-4=0$ ;
  \item $(d_{1})$ : $2x-3y+1=0$ et $(d_{3})$ : $4x-6y+5=0$ ;
  \item $(d_{4})$ : $x+5y-2=0$ et $(d_{5})$ : $y=2x+1$.
\end{questions}
\begin{correction}
\cas{a)}
\begin{calculs}
\vec{n}_{1} & \coord{2}{-3} &\qquad& \rem{normal à $(d_{1})$}\\
\vec{n}_{2} & \coord{3}{2} && \rem{normal à $(d_{2})$}\\
\vec{n}_{1}\cdot\vec{n}_{2} &= 2\times 3+(-3)\times 2 && \rem{$xx'+yy'$}\\
\vec{n}_{1}\cdot\vec{n}_{2} &= 0 &&
\end{calculs}
\explique{Les vecteurs normaux sont orthogonaux : $(d_{1})$ et $(d_{2})$ \res{sont
perpendiculaires}.}{propriété 11}
\cas{b)}
\begin{calculs}
\vec{n}_{3} & \coord{4}{-6} &\qquad& \rem{normal à $(d_{3})$}\\
\det(\vec{n}_{1}\,;\vec{n}_{3}) &= 2\times(-6)-(-3)\times 4 && \rem{$xy'-yx'$}\\
\det(\vec{n}_{1}\,;\vec{n}_{3}) &= 0 && \rem{$-12+12$}
\end{calculs}
\explique{Les vecteurs normaux sont colinéaires : $(d_{1})$ et $(d_{3})$ \res{sont
parallèles}.}{d'ailleurs $\vec{n}_{3}=2\,\vec{n}_{1}$}
\cas{c)}
\explique{On écrit $(d_{5})$ sous forme cartésienne : $2x-y+1=0$.}{on passe $y$ dans le
membre de droite}
\begin{calculs}
\vec{n}_{4} & \coord{1}{5} &\qquad& \rem{normal à $(d_{4})$}\\
\vec{n}_{5} & \coord{2}{-1} && \rem{normal à $(d_{5})$}\\
\vec{n}_{4}\cdot\vec{n}_{5} &= 1\times 2+5\times(-1) && \rem{$xx'+yy'$}\\
\vec{n}_{4}\cdot\vec{n}_{5} &= -3 &&\\
\det(\vec{n}_{4}\,;\vec{n}_{5}) &= 1\times(-1)-5\times 2 && \rem{$xy'-yx'$}\\
\det(\vec{n}_{4}\,;\vec{n}_{5}) &= -11 &&
\end{calculs}
\explique{Ni produit scalaire nul, ni déterminant nul : $(d_{4})$ et $(d_{5})$ \res{sont
sécantes, non perpendiculaires}.}{propriété 11}
\end{correction}

\etape{À vous de jouer}
Le repère est orthonormé. Pour chaque couple de droites, dire si elles sont parallèles,
perpendiculaires, ou ni l'un ni l'autre.
\begin{questions}[label=\alph*)]
  \item $(d_{1})$ : $5x+2y-1=0$ et $(d_{2})$ : $2x-5y+3=0$ ;
  \item $(d_{1})$ : $5x+2y-1=0$ et $(d_{3})$ : $10x+4y+7=0$ ;
  \item $(d_{4})$ : $3x-y+2=0$ et $(d_{5})$ : $y=x-4$.
\end{questions}
\reponse{a)~perpendiculaires ;\; b)~parallèles ;\; c)~sécantes, non perpendiculaires.}
\end{methodebox}

\afaire{exercices 34 à 37}

\partiecours{Projeté orthogonal d'un point sur une droite}

\begin{proprietebox}[Propriété 12 --- Coordonnées du projeté orthogonal]
\begin{minipage}[t]{0.58\linewidth}\raggedright
Le projeté orthogonal $H$ d'un point $M$ sur une droite $(d)$ (définition 5) est le point
d'intersection de $(d)$ et de la perpendiculaire $(\Delta)$ à $(d)$ qui passe par $M$.
\begin{itemize}
  \item Un vecteur \textbf{directeur} de $(d)$ est un vecteur \textbf{normal} à
        $(\Delta)$ : on trouve ainsi une équation de $(\Delta)$ (méthode 14).
  \item Les coordonnées de $H$ sont la solution du \textbf{système} formé par les
        équations de $(d)$ et de $(\Delta)$.
  \item La distance $MH$ est la \textbf{distance du point $M$ à la droite $(d)$} : c'est la
        plus courte distance de $M$ à un point de $(d)$.
\end{itemize}
\end{minipage}\hfill
\begin{minipage}[t]{0.38\linewidth}
\vspace{0pt}\centering
\begin{tikzpicture}[scale=0.6]
  \draw[fig] (-0.5,-0.25) -- (5,2.5) node[right,font=\footnotesize]{$(d)$};
  \draw[aide] (1.45,3.6) -- (2.75,1.0);
  \node[font=\footnotesize,left] at (1.5,3.5) {$(\Delta)$};
  \draw[black!60] (2.87,1.43) -- (2.74,1.70) -- (2.47,1.57);
  \pt{1.9,2.7}{right}{M} \pt{2.6,1.3}{below right}{H}
\end{tikzpicture}
\end{minipage}
\end{proprietebox}

\begin{remarquebox}[Rappel --- Système de deux équations à deux inconnues]
Le point d'intersection de deux droites sécantes a des coordonnées $(x\,;y)$ qui vérifient
les deux équations à la fois : le couple $(x\,;y)$ est l'unique solution du système
\[
  \left\{\begin{array}{l}
  2x+y-4=0\\
  -x+2y-3=0
  \end{array}\right.
  \qquad\text{(par exemple).}
\]
On le résout de deux façons :
\begin{itemize}
  \item par \textbf{substitution} : on isole une inconnue dans une équation, puis on la
        remplace dans l'autre ;
  \item par \textbf{combinaison} : on ajoute les deux équations, multipliées au besoin par
        des nombres bien choisis, pour faire disparaître une inconnue.
\end{itemize}
\end{remarquebox}

\begin{methodebox}[Méthode 16 --- Trouver le projeté orthogonal d'un point sur une droite]
Équation de la perpendiculaire, puis intersection des deux droites par un système.
\etapeexemples
Le repère est orthonormé. On considère la droite $(d)$ d'équation $2x+y-4=0$ et le point
$M(5\,;4)$.
\begin{questions}[label=\alph*)]
  \item Trouver une équation cartésienne de la perpendiculaire $(\Delta)$ à $(d)$ qui
        passe par $M$.
  \item Calculer les coordonnées du projeté orthogonal $H$ de $M$ sur $(d)$.
  \item En déduire la distance du point $M$ à la droite $(d)$.
  \item On considère les points $A(0\,;1)$, $B(4\,;3)$ et $C(1\,;4)$. Calculer les
        coordonnées du projeté orthogonal $K$ de $C$ sur la droite $(AB)$.
\end{questions}
\begin{correction}
\cas{a)}
\explique{$\vec{u}\coord{-1}{2}$ dirige $(d)$, donc il est normal à $(\Delta)$.}{$\vec{u}\coord{-b}{a}$
avec $a=2$ et $b=1$}
\begin{calculs}
-x+2y+c &= 0 &\qquad& \rem{$a=-1$ et $b=2$}\\
-5+2\times 4+c &= 0 && \rem{$M(5\,;4)$ est sur $(\Delta)$}\\
c &= -3 && \rem{$3+c=0$}
\end{calculs}
\explique{Une équation de $(\Delta)$ est $\res{-x+2y-3=0}$.}{on remplace $c$}
\cas{b)}
\explique{$H(x\,;y)$ est sur $(d)$ et sur $(\Delta)$ : $(x\,;y)$ est solution du système
$\left\{\begin{array}{l}2x+y-4=0\\-x+2y-3=0\end{array}\right.$}{équations de $(d)$, puis de
$(\Delta)$}
\explique{On le résout par substitution.}{$y$ s'isole facilement dans la première
équation}
\begin{calculs}
y &= 4-2x &\qquad& \rem{on isole $y$ dans la première équation}\\
-x+2(4-2x)-3 &= 0 && \rem{on remplace $y$ dans la seconde}\\
-x+8-4x-3 &= 0 && \rem{on développe}\\
-5x+5 &= 0 && \rem{on réduit}\\
x &= 1 &&\\
y &= 2 && \rem{$y=4-2\times 1$}
\end{calculs}
\explique{Le projeté orthogonal de $M$ sur $(d)$ est $\res{H(1\,;2)}$.}{on vérifie dans la
seconde équation : $-1+4-3=0$}
\cas{c)}
\begin{calculs}
\vect{MH} & \coord{-4}{-2} &\qquad& \rem{$1-5$ et $2-4$}\\
MH &= \sqrt{(-4)^{2}+(-2)^{2}} &&\\
MH &= \sqrt{20} &&\\
MH &= \res{2\sqrt{5}} && \rem{$\sqrt{4}\times\sqrt{5}$}
\end{calculs}
\cas{d)}
\explique{On cherche d'abord une équation de $(AB)$ : $\vect{AB}\coord{4}{2}$ la dirige,
donc $\vec{n}\coord{1}{-2}$ lui est normal.}{$\vect{AB}\cdot\vec{n}=4-4=0$}
\begin{calculs}
x-2y+c &= 0 &\qquad& \rem{équation de $(AB)$}\\
0-2\times 1+c &= 0 && \rem{$A(0\,;1)$ est sur $(AB)$}\\
c &= 2 &&
\end{calculs}
\explique{$(AB)$ : $x-2y+2=0$. La perpendiculaire à $(AB)$ qui passe par $C$ a pour
vecteur normal $\vect{AB}\coord{4}{2}$.}{puis la perpendiculaire}
\begin{calculs}
4x+2y+c &= 0 &\qquad& \rem{$a=4$ et $b=2$}\\
4\times 1+2\times 4+c &= 0 && \rem{$C(1\,;4)$ est sur la perpendiculaire}\\
c &= -12 &&
\end{calculs}
\explique{$K(x\,;y)$ est solution du système
$\left\{\begin{array}{l}x-2y+2=0\\4x+2y-12=0\end{array}\right.$}{équations de $(AB)$, puis de
la perpendiculaire}
\explique{On le résout par combinaison : en ajoutant les deux équations, $-2y$ et $+2y$
s'éliminent.}{les coefficients de $y$ sont opposés}
\begin{calculs}
5x-10 &= 0 &\qquad& \rem{$(x+4x)+(-2y+2y)+(2-12)$}\\
x &= 2 &&\\
2-2y+2 &= 0 && \rem{on remplace $x$ par $2$ dans la première équation}\\
y &= 2 && \rem{$4-2y=0$}
\end{calculs}
\explique{Le projeté orthogonal de $C$ sur $(AB)$ est $\res{K(2\,;2)}$.}{c'est le pied de la
hauteur issue de $C$ dans le triangle $ABC$}
\end{correction}

\etape{À vous de jouer}
Le repère est orthonormé. On considère la droite $(d)$ d'équation $3x-y-2=0$ et le point
$M(-2\,;2)$.
\begin{questions}[label=\alph*)]
  \item Trouver une équation cartésienne de la perpendiculaire $(\Delta)$ à $(d)$ qui
        passe par $M$.
  \item Calculer les coordonnées du projeté orthogonal $H$ de $M$ sur $(d)$.
  \item En déduire la distance du point $M$ à la droite $(d)$.
  \item On considère les points $E(1\,;0)$, $F(5\,;2)$ et $G(2\,;3)$. Calculer les
        coordonnées du projeté orthogonal $K$ de $G$ sur la droite $(EF)$.
\end{questions}
\reponse{a)~$x+3y-4=0$ ;\; b)~$H(1\,;1)$ ;\; c)~$\sqrt{10}$ ;\; d)~$K(3\,;1)$.}
\end{methodebox}

\afaire{exercices 38 à 40}

\partiecours{Calcul d'un angle}

\begin{proprietebox}[Propriété 13 --- Cosinus d'un angle]
Pour deux vecteurs non nuls $\vec{u}$ et $\vec{v}$, d'angle $\alpha$, la définition 4 donne
\[
  \cos(\alpha)=\frac{\vec{u}\cdot\vec{v}}{\norme{\vec{u}}\times\norme{\vec{v}}}.
\]
On calcule le produit scalaire autrement (coordonnées, normes) et on en déduit le
cosinus ; l'angle, entre $0^{\circ}$ et $180^{\circ}$, se lit dans le tableau des valeurs
remarquables ou s'obtient par la touche $\cos^{-1}$ de la calculatrice, en mode degré.
\end{proprietebox}

\begin{methodebox}[Méthode 17 --- Calculer un angle avec le produit scalaire]
\etapeexemples
Le repère est orthonormé.
\begin{questions}[label=\alph*)]
  \item Calculer la valeur exacte de l'angle $\alpha$ des vecteurs $\vec{u}\coord{3}{1}$ et
        $\vec{v}\coord{1}{2}$.
  \item On considère les points $A(1\,;1)$, $B(5\,;1)$ et $C(2\,;4)$. Calculer l'angle
        $\widehat{BAC}$, au dixième de degré.
  \item Dans le triangle $EFG$, $EF=5$, $EG=8$ et $FG=7$. Calculer l'angle $\widehat{FEG}$.
\end{questions}
\begin{correction}
\cas{a)}
\explique{On calcule le produit scalaire et les deux normes.}{les trois nombres de la
formule}
\begin{calculs}
\vec{u}\cdot\vec{v} &= 3\times 1+1\times 2 &\qquad& \rem{$xx'+yy'$}\\
\vec{u}\cdot\vec{v} &= 5 &&\\
\norme{\vec{u}} &= \sqrt{10} && \rem{$\sqrt{3^{2}+1^{2}}$}\\
\norme{\vec{v}} &= \sqrt{5} && \rem{$\sqrt{1^{2}+2^{2}}$}
\end{calculs}
\explique{On en déduit le cosinus, puis l'angle.}{propriété 13}
\begin{calculs}
\cos(\alpha) &= \frac{5}{\sqrt{10}\times\sqrt{5}} &\qquad& \rem{on remplace dans la formule}\\
\cos(\alpha) &= \frac{5}{\sqrt{50}} && \rem{$\sqrt{10}\times\sqrt{5}=\sqrt{50}$}\\
\cos(\alpha) &= \frac{5}{5\sqrt{2}} && \rem{$\sqrt{50}=\sqrt{25}\times\sqrt{2}$}\\
\cos(\alpha) &= \frac{\sqrt{2}}{2} && \rem{$\frac{1}{\sqrt{2}}=\frac{\sqrt{2}}{2}$}
\end{calculs}
\explique{L'angle mesure $\res{\alpha=45^{\circ}}$.}{valeur remarquable}
\cas{b)}
\begin{calculs}
\vect{AB} & \coord{4}{0} &\qquad& \rem{les deux vecteurs issus de $A$}\\
\vect{AC} & \coord{1}{3} &&\\
\vect{AB}\cdot\vect{AC} &= 4\times 1+0\times 3 && \rem{$xx'+yy'$}\\
\vect{AB}\cdot\vect{AC} &= 4 &&\\
AB &= 4 && \rem{$\sqrt{16+0}$}\\
AC &= \sqrt{10} && \rem{$\sqrt{1+9}$}\\
\cos\big(\widehat{BAC}\big) &= \frac{4}{4\times\sqrt{10}} && \rem{propriété 13}\\
\cos\big(\widehat{BAC}\big) &= \frac{1}{\sqrt{10}} &&\\
\widehat{BAC} &\approx \res{71{,}6^{\circ}} && \rem{touche $\cos^{-1}$, en mode degré}
\end{calculs}
\cas{c)}
\explique{Pas de coordonnées : on calcule le produit scalaire avec les normes.}{méthode
13, sommet commun $E$}
\begin{calculs}
\vect{EF}\cdot\vect{EG} &= \frac{1}{2}\left(EF^{2}+EG^{2}-FG^{2}\right) &\qquad& \rem{propriété 10}\\
\vect{EF}\cdot\vect{EG} &= \frac{1}{2}\left(25+64-49\right) &&\\
\vect{EF}\cdot\vect{EG} &= 20 &&\\
\cos\big(\widehat{FEG}\big) &= \frac{20}{5\times 8} && \rem{propriété 13}\\
\cos\big(\widehat{FEG}\big) &= \frac{1}{2} &&
\end{calculs}
\explique{L'angle mesure $\res{\widehat{FEG}=60^{\circ}}$.}{valeur remarquable}
\end{correction}

\etape{À vous de jouer}
Le repère est orthonormé.
\begin{questions}[label=\alph*)]
  \item Calculer la valeur exacte de l'angle $\alpha$ des vecteurs $\vec{u}\coord{3}{1}$ et
        $\vec{v}\coord{-1}{-2}$.
  \item On considère les points $E(0\,;1)$, $F(3\,;1)$ et $G(1\,;3)$. Calculer l'angle
        $\widehat{FEG}$, au dixième de degré.
  \item Dans le triangle $KLM$, $KL=3$, $KM=5$ et $LM=7$. Calculer l'angle $\widehat{LKM}$.
\end{questions}
\reponse{a)~$135^{\circ}$ ;\; b)~environ $63{,}4^{\circ}$ ;\; c)~$120^{\circ}$.}
\end{methodebox}

\afaire{exercices 41 à 43}

\partiecours{Théorème d'Al-Kashi}

\begin{proprietebox}[Théorème 6 --- Théorème d'Al-Kashi]
\begin{minipage}[t]{0.58\linewidth}\raggedright
Dans un triangle $ABC$, quand on connaît deux côtés et l'angle entre eux, on calcule le
troisième côté :
\[
  BC^{2}=AB^{2}+AC^{2}-2\times AB\times AC\times\cos\big(\widehat{BAC}\big).
\]
C'est le théorème de Pythagore « corrigé » : si l'angle en $A$ est droit,
$\cos(90^{\circ})=0$ et on retrouve $BC^{2}=AB^{2}+AC^{2}$.
\end{minipage}\hfill
\begin{minipage}[t]{0.38\linewidth}
\vspace{0pt}\centering
\begin{tikzpicture}[scale=0.6]
  \draw[fig] (0,0) -- (5,0) -- (1.5,2.6) -- cycle;
  \draw (0.7,0) arc (0:60:0.7);
  \pt{0,0}{below left}{A} \pt{5,0}{below right}{B} \pt{1.5,2.6}{above}{C}
\end{tikzpicture}
\end{minipage}

\smallskip
\demo{avec la relation de Chasles et l'identité remarquable (propriété 9) :}
\begin{calculs}
BC^{2} &= \norme{\vect{AC}-\vect{AB}}^{2} &\qquad& \rem{$\vect{BC}=\vect{BA}+\vect{AC}=\vect{AC}-\vect{AB}$}\\
BC^{2} &= AC^{2}-2\,\vect{AB}\cdot\vect{AC}+AB^{2} && \rem{identité remarquable}\\
BC^{2} &= AB^{2}+AC^{2}-2\times AB\times AC\times\cos\big(\widehat{BAC}\big) && \rem{définition 4}
\end{calculs}
\end{proprietebox}

\begin{methodebox}[Méthode 18 --- Calculer une longueur avec le théorème d'Al-Kashi]
On repère les deux côtés connus et l'angle compris entre eux ; le côté calculé est celui
qui fait face à cet angle.
\etapeexemples
\begin{questions}[label=\alph*)]
  \item Dans le triangle $ABC$, $AB=5$, $AC=8$ et $\widehat{BAC}=60^{\circ}$. Calculer $BC$.
  \item Deux voiliers quittent le même port. À midi, l'un est à $6$~km du port, l'autre à
        $10$~km, et leurs routes font un angle de $120^{\circ}$. Calculer la distance qui
        sépare les deux voiliers.
  \item Dans le triangle $DEF$, $DE=4$, $DF=3\sqrt{2}$ et $\widehat{EDF}=45^{\circ}$. Calculer
        la valeur exacte de $EF$.
\end{questions}
\begin{correction}
\cas{a)}
\begin{calculs}
BC^{2} &= AB^{2}+AC^{2}-2\times AB\times AC\times\cos\big(\widehat{BAC}\big) &\qquad& \rem{on recopie le théorème}\\
BC^{2} &= 25+64-2\times 5\times 8\times\frac{1}{2} && \rem{$\cos(60^{\circ})=\frac{1}{2}$}\\
BC^{2} &= 49 && \rem{$89-40$}\\
BC &= \res{7} && \rem{une longueur est positive : $\sqrt{49}$}
\end{calculs}
\cas{b)}
\explique{Le port $O$ et les voiliers $V_{1}$ et $V_{2}$ forment un triangle, avec $OV_{1}=6$,
$OV_{2}=10$ et $\widehat{V_{1}OV_{2}}=120^{\circ}$.}{on nomme les points, puis on repère
l'angle}
\begin{calculs}
V_{1}V_{2}^{2} &= 6^{2}+10^{2}-2\times 6\times 10\times\cos(120^{\circ}) &\qquad& \rem{théorème d'Al-Kashi}\\
V_{1}V_{2}^{2} &= 36+100-120\times\left(-\frac{1}{2}\right) && \rem{$\cos(120^{\circ})=-\frac{1}{2}$}\\
V_{1}V_{2}^{2} &= 196 && \rem{$136+60$ : moins par moins}\\
V_{1}V_{2} &= 14 && \rem{$\sqrt{196}$}
\end{calculs}
\explique{Les deux voiliers sont à $\res{14~\text{km}}$ l'un de l'autre.}{on conclut avec
l'unité}
\cas{c)}
\begin{calculs}
EF^{2} &= DE^{2}+DF^{2}-2\times DE\times DF\times\cos\big(\widehat{EDF}\big) &\qquad& \rem{angle en $D$, entre $[DE]$ et $[DF]$}\\
EF^{2} &= 16+18-2\times 4\times 3\sqrt{2}\times\frac{\sqrt{2}}{2} && \rem{$\big(3\sqrt{2}\big)^{2}=9\times 2$}\\
EF^{2} &= 34-24 && \rem{$\sqrt{2}\times\sqrt{2}=2$}\\
EF^{2} &= 10 &&\\
EF &= \res{\sqrt{10}} &&
\end{calculs}
\end{correction}

\etape{À vous de jouer}
\begin{questions}[label=\alph*),beginpenalty=10000]
  \item Dans le triangle $KLM$, $KL=8$, $KM=15$ et $\widehat{LKM}=60^{\circ}$. Calculer $LM$.
  \item Deux randonneurs partent du même refuge. Une heure plus tard, l'un est à $3$~km du
        refuge, l'autre à $5$~km, et leurs chemins font un angle de $120^{\circ}$.
        Calculer la distance qui sépare les deux randonneurs.
  \item Dans le triangle $RST$, $RS=2\sqrt{2}$, $RT=5$ et $\widehat{SRT}=45^{\circ}$. Calculer
        la valeur exacte de $ST$.
\end{questions}
\reponse{a)~$LM=13$ ;\; b)~$7$~km ;\; c)~$ST=\sqrt{13}$.}
\end{methodebox}

\afaire{exercices 44 à 47}

\partiecours{Ensemble des points \texorpdfstring{$M$ tels que $\vect{MA}\cdot\vect{MB}=0$}{M tels que MA.MB = 0}}

\begin{proprietebox}[Théorème 7 --- Cercle de diamètre {$[AB]$}]
\begin{minipage}[t]{0.58\linewidth}\raggedright
Soit $A$ et $B$ deux points distincts. L'ensemble des points $M$ tels que
$\vect{MA}\cdot\vect{MB}=0$ est le \textbf{cercle de diamètre $[AB]$}.

\smallskip
Pour $M$ distinct de $A$ et de $B$, cela veut dire : le triangle $AMB$ est rectangle en
$M$ si, et seulement si, $M$ est sur le cercle de diamètre $[AB]$.
\end{minipage}\hfill
\begin{minipage}[t]{0.38\linewidth}
\vspace{0pt}\centering
\begin{tikzpicture}[scale=0.55]
  \draw[coulCourbe,line width=1pt] (0,0) circle (2);
  \draw[fig] (-2,0) -- (1.2,1.6) -- (2,0);
  \draw[aide] (-2,0) -- (2,0);
  \draw[black!60] (0.93,1.47) -- (1.08,1.2) -- (1.35,1.33);
  \pt{-2,0}{left}{A} \pt{2,0}{right}{B} \pt{0,0}{below}{I} \pt{1.2,1.6}{above right}{M}
\end{tikzpicture}
\end{minipage}

\smallskip
\demo{avec le milieu $I$ de $[AB]$ : $\vect{IB}=-\vect{IA}$, et $IA=\frac{AB}{2}$.}
\begin{calculs}
\vect{MA}\cdot\vect{MB} &= \big(\vect{MI}+\vect{IA}\big)\cdot\big(\vect{MI}+\vect{IB}\big) &\qquad& \rem{relation de Chasles}\\
\vect{MA}\cdot\vect{MB} &= \big(\vect{MI}+\vect{IA}\big)\cdot\big(\vect{MI}-\vect{IA}\big) && \rem{$\vect{IB}=-\vect{IA}$}\\
\vect{MA}\cdot\vect{MB} &= MI^{2}-IA^{2} && \rem{identité remarquable (propriété 9)}
\end{calculs}
Donc $\vect{MA}\cdot\vect{MB}=0$ quand $MI^{2}=IA^{2}$, c'est-à-dire $MI=IA$ (des
distances sont positives) : $M$ est sur le cercle de centre $I$ et de rayon $IA$.
\end{proprietebox}

\begin{methodebox}[Méthode 19 --- Trouver un ensemble de points avec le produit scalaire]
On traduit la condition par le produit scalaire, puis on reconnaît un cercle
(théorème 7) ou une droite (théorème 5).
\etapeexemples
Le repère est orthonormé.
\begin{questions}[label=\alph*)]
  \item On considère les points $A(-1\,;2)$ et $B(3\,;0)$. Trouver une équation de
        l'ensemble $(\mathcal{C})$ des points $M$ tels que $\vect{MA}\cdot\vect{MB}=0$.
  \item Le point $C(0\,;-1)$ est-il sur $(\mathcal{C})$ ? En déduire la nature du triangle
        $ABC$.
  \item On considère les points $D(1\,;2)$ et $E(4\,;-1)$. Trouver l'ensemble $(\mathcal{D})$
        des points $M$ tels que $\vect{DE}\cdot\vect{DM}=0$, et une équation de cet ensemble.
\end{questions}
\begin{correction}
\cas{a)}
\explique{$(\mathcal{C})$ est le cercle de diamètre $[AB]$ : son centre est le milieu $I$ de
$[AB]$, son rayon est $IA$.}{théorème 7}
\begin{calculs}
I & \left(\dfrac{-1+3}{2}\,;\dfrac{2+0}{2}\right) &\qquad& \rem{coordonnées du milieu}\\
I & (1\,;1) &&\\
IA^{2} &= (-1-1)^{2}+(2-1)^{2} && \rem{le carré du rayon suffit}\\
IA^{2} &= 5 &&
\end{calculs}
\explique{Une équation de $(\mathcal{C})$ est $\res{(x-1)^{2}+(y-1)^{2}=5}$.}{théorème 2 :
$(x-a)^{2}+(y-b)^{2}=r^{2}$}
\cas{b)}
\begin{calculs}
(0-1)^{2}+(-1-1)^{2} &= 1+4 &\qquad& \rem{on remplace $x$ par $0$ et $y$ par $-1$}\\
(0-1)^{2}+(-1-1)^{2} &= 5 &&
\end{calculs}
\explique{L'équation est vérifiée : $C$ est sur $(\mathcal{C})$, donc
$\vect{CA}\cdot\vect{CB}=0$ et le triangle $ABC$ \res{est rectangle en $C$}.}{théorème 7}
\cas{c)}
\explique{$\vect{DM}$ est orthogonal à $\vect{DE}$ : $(\mathcal{D})$ est \res{la droite qui
passe par $D$, perpendiculaire à $(DE)$}, de vecteur normal $\vect{DE}$.}{théorème 5}
\begin{calculs}
\vect{DE} & \coord{3}{-3} &\qquad& \rem{$4-1$ et $-1-2$}\\
\vect{DM} & \coord{x-1}{y-2} && \rem{avec $M(x\,;y)$}\\
\vect{DE}\cdot\vect{DM} &= 3(x-1)-3(y-2) && \rem{$xx'+yy'$}\\
\vect{DE}\cdot\vect{DM} &= 3x-3y+3 && \rem{$3x-3-3y+6$}
\end{calculs}
\explique{Une équation de $(\mathcal{D})$ est $\res{3x-3y+3=0}$, ou $x-y+1=0$.}{on divise
par $3$}
\end{correction}

\etape{À vous de jouer}
Le repère est orthonormé.
\begin{questions}[label=\alph*)]
  \item On considère les points $F(-2\,;1)$ et $G(4\,;-1)$. Trouver une équation de
        l'ensemble $(\mathcal{C})$ des points $M$ tels que $\vect{MF}\cdot\vect{MG}=0$.
  \item Le point $K(2\,;3)$ est-il sur $(\mathcal{C})$ ? En déduire la nature du triangle
        $FGK$.
  \item On considère les points $R(2\,;-1)$ et $S(0\,;3)$. Trouver l'ensemble $(\mathcal{D})$
        des points $M$ tels que $\vect{RS}\cdot\vect{RM}=0$, et une équation de cet ensemble.
\end{questions}
\reponse{a)~$(x-1)^{2}+y^{2}=10$ ;\; b)~oui, $FGK$ est rectangle en $K$ ;\; c)~la
perpendiculaire à $(RS)$ en $R$, d'équation $x-2y-4=0$.}
\end{methodebox}

\afaire{exercices 48 à 50}

\end{document}
